% Corrections to "Theory of the Electron-Phonon-Photon coupling" (draft v2).
% Standalone document; packages amsmath, amssymb, bm only; no custom macros in any display.
% This file supersedes corrections/issue-phase.tex, issue-e410.tex and issue-efactor.tex.
% Supporting derivations: corrections/issue-*-report.md. Numerical checks (git-ignored):
% test/corr-phase/check_phase.py, test/corr-e410/check_e410.py, test/corr-efactor/check_efactor.py,
% test/corr-merge/check_merge.py (merged forms of Sec. 5).
\documentclass[11pt]{article}
\usepackage{amsmath,amssymb,bm}
\setlength{\textwidth}{6.5in}
\setlength{\oddsidemargin}{0in}
\setlength{\evensidemargin}{0in}
\setlength{\textheight}{8.9in}
\setlength{\topmargin}{-0.4in}
\allowdisplaybreaks

\title{Corrections to `Theory of the Electron-Phonon-Photon coupling' (draft v2)}
\author{}
\date{25 October 2026}

\begin{document}
\maketitle

\section{Introduction}\label{sec:intro}

\paragraph{What was checked, and how.}
While an interactive primer on F.~Giustino, Rev.\ Mod.\ Phys.\ \textbf{89}, 015003 (2017) and on the
present manuscript was being written, every displayed equation of draft v2 (102 pages) was checked
against the page images of the PDF. Three issues go beyond typos. They change equations, or the reading
of a displayed result:
\begin{enumerate}
\item The Born-charge dipoles are placed at the cell origin $\mathbf R^0_i$ instead of the atomic site
$\mathbf R^0_{ia}=\mathbf R^0_i+\mathbf r^0_{\rm a}$ (Sec.~3 and everything built on it). The basis phase
$e^{-i(\mathbf q+\mathbf G)\mathbf r^0_{\rm a}}$ is restored in Sec.~\ref{sec:phase}.
\item Eq.~(410) is the $\mathbf q\to0$, finite-$c$, electron-free end point of the upper phonon
polariton, not the LO--TO splitting (Sec.~7.1). It differs from the electrostatic result (454) in two
independent ways, electronic screening and the order of limits (Sec.~\ref{sec:e410}).
\item The charge $e$ is lost in the continuity equation on p.~94, so the electrostatic Fr\"ohlich
vertices on pp.~95 and 98 lack a factor $e$ (Sec.~7.2); Sec.~\ref{sec:efactor}.
\end{enumerate}
Each correction was derived from the manuscript's own definitions, checked by two independent
skeptical reviewers working from the page images, and revised until every objection was answered.
Each was tested numerically: dimensional bookkeeping, limits, model calculations for GaAs and for
random multi-atom cells, and comparison with Eq.~(180) of Giustino's review. Two corrections touch
the same equations on pp.~95--99, and the phase also enters the general-$\mathbf q$ result of
Sec.~\ref{sec:e410}. The merged forms were derived again and checked numerically; they are
collected in Sec.~\ref{sec:interplay}, which also fixes the order of the new material on p.~99.
Sec.~\ref{sec:minor} lists the smaller errata found during the check.

\paragraph{Conventions.}
Page numbers are those of the PDF. ``Eq.~($n$)'' is equation ($n$) of the draft. The manuscript's
notation is used throughout: roman a, b label the atoms of a cell, $ia$ the atom a of cell $i$,
$\Omega=N\Omega_{\rm uc}$, and $\sigma$ the phonon branches from Sec.~5 on ($\alpha$ in Sec.~3).
\begin{itemize}
\item \emph{Units.} Sec.~3 is in SI units, as stated on p.~23, and writes the electron charge as $-e$
with $e>0$ [before (123)]. From Sec.~4 on the units are Gaussian ($e^2/4\pi\epsilon_0\to e^2$,
$1/\epsilon_0\to4\pi$), and $e$ is the signed electron charge, $e=-|e|$, as in the Lagrangian (134). All
corrected equations of Secs.~4--8 are Gaussian.
\item \emph{Born charges} $Z^*_{{\rm a}\mu\nu}$ have the dimension of charge [Eq.~(110)]:
$Z^*_{{\rm a}\mu\nu}=|e|\,Z^*_{\kappa,\mu\nu}$, with Giustino's dimensionless $Z^*_\kappa$. In Gaussian units
$eZ^*$ has the dimension of energy$\times$length.
\item \emph{Densities.} $\rho_{\rm e}=\Psi^+\Psi$ [Eqs.~(239), (420)] is the electron \emph{number}
density and $\rho_{n\mathbf k,m\mathbf k+\mathbf q}$ [Eqs.~(106), (426)] its dimensionless matrix element.
Wherever a charge density is meant, $e\rho_{\rm e}$ is written (Sec.~\ref{sec:efactor}).
\item \emph{Phonon phase.} Eqs.~(7), (83) and (242) attach the Bloch phase to the cell,
$Q_{ia\nu}\propto e_{{\rm a}\nu\sigma}(\mathbf q)\,e^{i\mathbf q\mathbf R^0_i}$, as Giustino does. The
$e_{{\rm a}\nu\sigma}(\mathbf q)$ are the eigenvectors that solve the equations of motion for this pattern
(see Sec.~\ref{sec:phase-problem} on the sign in Eq.~(15)).
\end{itemize}

\paragraph{Equation tags.} One convention is used for every corrected and every new equation.
\begin{itemize}
\item ($n'$) is the corrected form of the numbered equation ($n$) and replaces it. A combined tag such
as (241$'$, 314$'$) means that one display replaces two printed equations with the same content; paste
it at each place, with the single tag.
\item (p$N'$), (p$N$a$'$), (p$N$b$'$), \ldots\ is the corrected form of the unnumbered display on page
$N$ (the first, second, \ldots\ such display on that page, as stated in the text) and replaces it.
\item ($n$a), ($n$b), \ldots\ is a new equation, inserted after Eq.~($n$) or after the passage that
follows it, as stated.
\item A slip of a single symbol inside a numbered equation is corrected in place, and the equation keeps
its number (example: $\Omega^0_\sigma\to\Omega^0_{\sigma\sigma}$ in Eq.~(406)).
\item Untagged displays are explanations and are not meant to be pasted.
\end{itemize}
\emph{LaTeX.} Every tagged display is an \verb|equation*| or \verb|align*| environment (with
\verb|split|, \verb|aligned| or \verb|gathered| inside where needed) and carries an explicit
\verb|\tag|. These environments never step the equation counter, and this document contains no
\verb|\stepcounter|.

\paragraph{How to paste.}
The displays need only \texttt{amsmath}, \texttt{amssymb} and \texttt{bm}, and no macros of their own.
\begin{itemize}
\item \emph{A primed equation} replaces a numbered \verb|equation|. Delete the old environment, paste
the \verb|equation*| block, and put \verb|\stepcounter{equation}| directly after it, so that all later
equations keep their printed numbers. Do not add \verb|\tag| to an unstarred \verb|equation| instead:
amsmath steps the counter there even when \verb|\tag| is given, and the two methods must not be mixed.
\item \emph{A page-tagged display} replaces an unnumbered display, and \emph{a lettered equation} is new.
Paste them as they are; no counter needs to be stepped.
\item If the next version is renumbered from scratch, drop all \verb|\tag|s, use the unstarred
environments, and keep the tags of this document only as a key to the changes.
\item Replacement prose is set in \texttt{quote} environments. ``Replace `\ldots' by \ldots'' quotes the
printed text. Where a correction is a rule applied to many displays, the affected displays are listed.
\end{itemize}

%======================================================================================
\section{The basis phase: Born dipoles sit at $\mathbf R^0_i+\mathbf r^0_{\rm a}$ (Sec.~3 and downstream)}\label{sec:phase}

\subsection{Location}
Sec.~3, pp.~24--25: Eqs.~(121)--(124) and (131). The same placement recurs in Eqs.~(134),
(141)--(143), (167), (183), (186) and the Hamiltonians between (184) and (186), (237), (241),
(270)--(272), (274), (286)--(311) and the displays between them, (314)--(315), and (352)--(354) with
the displays on pp.~76--77. Through the ionic vertex $j_{\alpha\sigma}(\mathbf q)$ [Eq.~(295) and the two
displays after (385)] it reaches every equation of Secs.~6.1--8. The vertex is written out explicitly in
the two displays after (417) on p.~93, in (418), (430), (447) and (448), and in the displays on pp.~95,
96 and 98.

\subsection{Problem}\label{sec:phase-problem}
Eqs.~(121)--(123) place the dipole $\mathbf p_{ia}$ of atom a at the cell origin $\mathbf R^0_i$. The atom
sits at $\mathbf R^0_{ia}=\mathbf R^0_i+\mathbf r^0_{\rm a}$, Eq.~(94). Sec.~2 keeps the corresponding factor,
$V_{\rm a}(\mathbf q')=e^{-i\mathbf q'\mathbf r^0_{\rm a}}v_{\rm a}(\mathbf q')$ in Eq.~(101), but
$V^{\mathcal L}_{\rm a}$ in Eq.~(124) does not. As a result Eq.~(131) lacks the factor
$e^{-i(\mathbf q+\mathbf G)\mathbf r^0_{\rm a}}$ inside the atom sum. Eq.~(180) of Giustino's review carries
it, and Giustino's footnote~7 notes that Eq.~(4) of C.~Verdi and F.~Giustino, Phys.\ Rev.\ Lett.\
\textbf{115}, 176401 (2015) lacked it too.

\emph{Why the phase is needed.} Eqs.~(7), (83) and (242) attach the Bloch phase to the cell:
$Q_{ia\nu}\propto e_{{\rm a}\nu\alpha}(\mathbf q)\,e^{i\mathbf q\mathbf R^0_i}$. The dipole
$Z^*_{{\rm a}\mu\nu}Q_{ia\nu}$ located at $\mathbf R^0_i+\mathbf r^0_{\rm a}$ therefore has, at wave vector
$\mathbf q+\mathbf G$, the amplitude $\propto Z^*_{{\rm a}\mu\nu}e_{{\rm a}\nu\alpha}(\mathbf q)\,
e^{-i(\mathbf q+\mathbf G)\mathbf r^0_{\rm a}}$, and the basis phase must appear explicitly in every vertex.
The sign of the phase follows from (83) and the Fourier convention (117) alone. Only the product
$e_{{\rm a}\nu\alpha}(\mathbf q)\,e^{-i(\mathbf q+\mathbf G)\mathbf r^0_{\rm a}}$ is independent of how the
atoms are grouped into cells. Moving atom a to a neighbouring cell gives
\begin{equation*}
\mathbf r^0_{\rm a}\to\mathbf r^0_{\rm a}+\mathbf R_L,\qquad
e_{{\rm a}\nu\alpha}(\mathbf q)\to e^{i\mathbf q\mathbf R_L}e_{{\rm a}\nu\alpha}(\mathbf q),\qquad
e_{{\rm a}\nu\alpha}(\mathbf q)\,e^{-i(\mathbf q+\mathbf G)\mathbf r^0_{\rm a}}\ \text{unchanged}.
\tag{131a}
\end{equation*}
With atom phases instead, $Q_{ia\nu}\propto\tilde e_{{\rm a}\nu\alpha}(\mathbf q)e^{i\mathbf q\mathbf R^0_{ia}}$, the
factor reduces to $e^{-i\mathbf G\mathbf r^0_{\rm a}}$. Eq.~(131) as printed is therefore correct only in
that convention and only at $\mathbf G=0$.

\emph{A precondition: the sign in Eqs.~(15) and (35).} The transformation law (131a) and the
cancellation (131b) below assume that the $e_{{\rm a}\nu\alpha}(\mathbf q)$ solve the equations of motion
for the pattern $e^{i\mathbf q\mathbf R^0_i}$, as in Giustino's Eqs.~(14) and (20). They are then the
eigenvectors of $(m_{\rm a}m_{\rm b})^{-1/2}\sum_{\bm{\delta}}e^{-i\mathbf q\bm{\delta}}
k_{{\rm a}\mu{\rm b}\nu}(\bm{\delta})$. Eqs.~(15) and (35) as printed have $e^{+i\mathbf q\bm{\delta}}$, and the
step after Eq.~(43) uses $\tilde\Omega^2_{{\rm b}\nu{\rm a}\mu}=\tilde\Omega^2_{{\rm a}\mu{\rm b}\nu}$, which
holds only for real $\tilde\Omega^2$. The eigenvectors of the printed matrix are the complex conjugates.
For a long acoustic wave they carry $e^{-i\mathbf q\mathbf r^0_{\rm a}}$, and the basis phase would then
double the spurious acoustic term of (131b) instead of removing it. The sign in (15) and (35) must
therefore be corrected together with the present issue (Sec.~\ref{sec:minor}, p.~4).

\subsection{Corrected equations, Sec.~3 (SI units, as in the manuscript)}
On p.~24, before (121), read: ``\ldots\ and recalling that the dipole of atom a in the $i$th cell sits at
the atomic site $\mathbf R^0_{ia}=\mathbf R^0_i+\mathbf r^0_{\rm a}$ [Eq.~(94)], we conclude that the potential
produced by the displacement $Q_{ia\nu}$ of atom a in the $\nu$ direction is''. The left-hand label of
(121) becomes $\nu$, the direction of the displacement.
\begin{align*}
\phi_{ia\nu}(\mathbf r)&=-i\sum_\mu\frac{Z^*_{{\rm a}\mu\nu}Q_{ia\nu}}{\epsilon_0}\int\frac{d\mathbf q}{(2\pi)^3}\,
e^{i\mathbf q(\mathbf r-\mathbf R^0_i-\mathbf r^0_{\rm a})}\,\frac{q_\mu}{\mathbf q\hat\epsilon^\infty\mathbf q},\tag{121$'$}\\
\phi_i(\mathbf r)&=-i\sum_{{\rm a}\mu\nu}\frac{Z^*_{{\rm a}\mu\nu}Q_{ia\nu}}{\epsilon_0}\int\frac{d\mathbf q}{(2\pi)^3}\,
e^{i\mathbf q(\mathbf r-\mathbf R^0_i-\mathbf r^0_{\rm a})}\,\frac{q_\mu}{\mathbf q\hat\epsilon^\infty\mathbf q},\tag{122$'$}\\
\Delta V^{\mathcal L}_{\rm cr}(\mathbf r)&=i\sum_{ia\mu\nu}Z^*_{{\rm a}\mu\nu}Q_{ia\nu}\int\frac{d\mathbf q}{(2\pi)^3}\,
\frac{e\,q_\mu}{\epsilon_0\,\mathbf q\hat\epsilon^\infty\mathbf q}\,e^{i\mathbf q(\mathbf r-\mathbf R^0_i-\mathbf r^0_{\rm a})},\tag{123$'$}\\
V^{\mathcal L}_{\rm a}(\mathbf q)&=e^{-i\mathbf q\mathbf r^0_{\rm a}}\,\frac{e^2}{\epsilon_0\,\mathbf q\hat\epsilon^\infty\mathbf q}.\tag{124$'$}
\end{align*}
Eqs.~(125)--(130) keep their form. With (124$'$) the identity between (128) and (102) now includes the
basis phase of (101), as it should. Inserting (124$'$) into (130) gives
\begin{equation*}
\begin{split}
g^{\mathcal L}_{nm\alpha}(\mathbf k,\mathbf q)={}&\frac{4\pi i}{\Omega_{\rm uc}}\frac{e}{4\pi\epsilon_0}
\sum_{\mathbf G}\sum_{{\rm a}\mu\nu}\sqrt{\frac{\hbar}{2Nm_{\rm a}\omega_\alpha(\mathbf q)}}\\
&\times\frac{(q_\mu+G_\mu)Z^*_{{\rm a}\mu\nu}e_{{\rm a}\nu\alpha}(\mathbf q)\,e^{-i(\mathbf q+\mathbf G)\mathbf r^0_{\rm a}}}
{(\mathbf q+\mathbf G)\hat\epsilon^\infty(\mathbf q+\mathbf G)}\,\rho_{n\mathbf k,m\mathbf k-\mathbf q}(-\mathbf G),
\end{split}
\tag{131$'$}
\end{equation*}
with $e/4\pi\epsilon_0\to e$ in Gaussian units. For an acoustic branch, with $M=\sum_{\rm a}m_{\rm a}$,
$\mathbf u$ the common displacement and $\mathbf w_{\rm a}(\mathbf q)=O(q)$ the internal-strain relaxation
(which does not depend on the choice of cell),
\begin{equation*}
e_{{\rm a}\nu\alpha}(\mathbf q)=e^{i\mathbf q\mathbf r^0_{\rm a}}\Big[\sqrt{m_{\rm a}/M}\,u_\nu+i\,w_{{\rm a}\nu}(\mathbf q)+O(q^2)\Big].
\end{equation*}
Then
\begin{equation*}
\begin{aligned}
\sum_{\rm a}\frac{\hat Z^*_{\rm a}\mathbf e_{{\rm a}\alpha}(\mathbf q)}{\sqrt{m_{\rm a}}}
-\sum_{\rm a}\frac{\hat Z^*_{\rm a}\mathbf e_{{\rm a}\alpha}(\mathbf q)}{\sqrt{m_{\rm a}}}\,e^{-i\mathbf q\mathbf r^0_{\rm a}}
&=\frac{i}{\sqrt M}\sum_{\rm a}(\mathbf q\mathbf r^0_{\rm a})\,\hat Z^*_{\rm a}\mathbf u+O(q^2),\\
\sum_{\rm a}\frac{\hat Z^*_{\rm a}\mathbf e_{{\rm a}\alpha}(\mathbf q)}{\sqrt{m_{\rm a}}}\,e^{-i\mathbf q\mathbf r^0_{\rm a}}
&=\frac{1}{\sqrt M}\Big(\sum_{\rm a}\hat Z^*_{\rm a}\Big)\mathbf u
+i\sum_{\rm a}\frac{\hat Z^*_{\rm a}\mathbf w_{\rm a}(\mathbf q)}{\sqrt{m_{\rm a}}}+O(q^2)\\
&=i\sum_{\rm a}\frac{\hat Z^*_{\rm a}\mathbf w_{\rm a}(\mathbf q)}{\sqrt{m_{\rm a}}}+O(q^2).
\end{aligned}
\tag{131b}
\end{equation*}
The printed (phase-less) vertex therefore equals the corrected one plus the spurious term
$(i/\sqrt M)\sum_{\rm a}(\mathbf q\mathbf r^0_{\rm a})\hat Z^*_{\rm a}\mathbf u$. Because $\omega_\alpha\propto q$,
this term gives $|g^{\mathcal L}|^2\propto1/q$ for acoustic phonons, the scaling of a piezoelectric
coupling. It depends on the choice of cell, and it survives in centrosymmetric crystals, where the
physical term $\sum_{\rm a}\hat Z^*_{\rm a}\mathbf w_{\rm a}$ vanishes. For GaAs the spurious term alone
corresponds to a piezoelectric constant $|\sum_{\rm a}Z^*_{\rm a}(\hat{\mathbf q}\mathbf r^0_{\rm a})|/\Omega_{\rm uc}$
of 0.6--1.8~C/m$^2$ along [111], depending on the cell, and of 1.04~C/m$^2$ along [100]. This is
comparable to or larger than the internal-strain coupling of the corrected vertex: about 1.2~C/m$^2$
along [111] in a central-force model of GaAs (nearest- and next-nearest-neighbour springs fitted to
$\omega_{\rm TO}$ and to the [100] sound velocity), and zero along [100]. The printed total can be larger
or smaller than the corrected one, depending on the cell (1.8 or 0.6~C/m$^2$ along [111] in that model).
It is also nonzero along [100], where the true coupling vanishes by symmetry. With the phase [second
line of (131b)] the leading term vanishes by the acoustic sum rule $\sum_{\rm a}\hat Z^*_{\rm a}=0$, and the
remainder is the internal-strain dipole. The measured coupling, $(2/\sqrt3)|e_{14}|\approx0.19$~C/m$^2$
for LA[111], is much smaller still: the clamped-ion (electronic) piezoelectric response largely cancels
the internal-strain term [S.~de~Gironcoli, S.~Baroni and R.~Resta, Phys.\ Rev.\ Lett.\ \textbf{62}, 2853
(1989)]. In the long-range vertex that response enters through the dynamical quadrupoles, so a
consistent $O(q)$ theory needs them [G.~Brunin \textit{et al.}, Phys.\ Rev.\ Lett.\ \textbf{125}, 136601
(2020); V.~A.~Jhalani \textit{et al.}, \textit{ibid.}\ 136602 (2020); J.~Park \textit{et al.}, Phys.\ Rev.\ B
\textbf{102}, 125203 (2020)]. These references are quoted from memory and should be verified.

\subsection{Corrected equations, Secs.~4--5 (Gaussian units)}
\emph{Rule.} In Secs.~4--6, in every display, numbered or not, replace $A_\alpha(\mathbf R^0_i)$,
$\dot A_\nu(\mathbf R^0_i)$ and $E_\nu(\mathbf R^0_i)$ by $A_\alpha(\mathbf R^0_{ia})$, $\dot A_\nu(\mathbf R^0_{ia})$ and
$E_\nu(\mathbf R^0_{ia})$, and replace $\delta(\mathbf r-\mathbf R^0_i)$, $\delta(\mathbf r'-\mathbf R^0_j)$ by
$\delta(\mathbf r-\mathbf R^0_{ia})$, $\delta(\mathbf r'-\mathbf R^0_{jb})$. Here a (b) is the atom index already
summed in the same term. The rule covers (134), (141)--(143), (167), (183), the Hamiltonian after (184)
on p.~38, the two Hamiltonians before (186) on p.~39, (186), (237), (241), (270)--(272), (286), the
display after (288) on p.~61, (289), (290), (302), (314), (315) and the two displays on p.~76 after (353).
In the Fourier substitutions (291), (303) and (354), and in the 2D substitutions on pp.~64 and 66, use
$\delta(\mathbf r-\mathbf R^0_{ia})\to\int\frac{d\mathbf q_1}{(2\pi)^3}e^{i\mathbf q_1(\mathbf r-\mathbf R^0_i-\mathbf r^0_{\rm a})}$,
or in 2D $\delta(z-z_{\rm a})\int\frac{d\mathbf Q_1}{(2\pi)^2}
e^{i\mathbf Q_1(\bm{\rho}-\mathbf R^0_i-\bm{\rho}^0_{\rm a})}$. In (287)--(290) the factor
$e^{i\mathbf q(\mathbf r-\mathbf r')}$ that comes from $\pi_{{\rm a}\mu}(\mathbf r)$ becomes
$e^{i\mathbf q(\mathbf R^0_i-\mathbf R^0_j)}$ by (274$'$). The displayed results are
\begin{align*}
L&=\dots+\frac1c\sum_{ia}\dot{\mathbf p}_{ia}\mathbf A(\mathbf R^0_{ia}),\tag{134$'$}\\
m_{\rm a}\ddot Q_{ia\mu}+\sum_{jb\nu}k_{ia\mu,jb\nu}Q_{jb\nu}&=\sum_\nu Z^*_{{\rm a}\nu\mu}E_\nu(\mathbf R^0_{ia}),\tag{143$'$}\\
\mathbf P(\mathbf r)&=\sum_{ia}\mathbf p_{ia}\,\delta(\mathbf r-\mathbf R^0_{ia}),\tag{167$'$}\\
\pi_{ia\mu}&=m_{\rm a}\dot Q_{ia\mu}+\frac1c\sum_\alpha Z^*_{{\rm a}\alpha\mu}A_\alpha(\mathbf R^0_{ia}),\tag{183$'$}
\end{align*}
\begin{equation*}
\begin{split}
V={}&-\frac{e}{2m_{\rm e}c}\sum_e\big[\mathbf p_e\mathbf A(\mathbf r_e)+\mathbf A(\mathbf r_e)\mathbf p_e\big]
+\frac{e^2}{2m_{\rm e}c^2}\sum_e\mathbf A^2(\mathbf r_e)\\
&-\frac1c\sum_{ia\alpha\mu}\frac{1}{m_{\rm a}}Z^*_{{\rm a}\alpha\mu}\pi_{ia\mu}A_\alpha(\mathbf R^0_{ia})
+\frac{1}{c^2}\sum_{ia\mu\alpha\beta}\frac{1}{2m_{\rm a}}Z^*_{{\rm a}\alpha\mu}Z^*_{{\rm a}\beta\mu}
A_\alpha(\mathbf R^0_{ia})A_\beta(\mathbf R^0_{ia}),
\end{split}
\tag{186$'$}
\end{equation*}
\begin{align*}
P_\mu(\mathbf r)&=\sum_{ia\nu}Z^*_{{\rm a}\mu\nu}Q_{ia\nu}\,\delta(\mathbf r-\mathbf R^0_{ia}),\tag{241$'$, 314$'$}\\
\rho^{\rm ion}_{\alpha\beta}(\mathbf r)&=\frac{m_0}{e^2}\sum_{ia\mu}\frac{1}{m_{\rm a}}Z^*_{{\rm a}\alpha\mu}
Z^*_{{\rm a}\beta\mu}\,\delta(\mathbf r-\mathbf R^0_{ia}),\tag{271$'$, 315$'$}\\
j^{\rm ion}_\alpha(\mathbf r)&=\sum_{ia\mu}\frac{1}{m_{\rm a}}Z^*_{{\rm a}\alpha\mu}\pi_{{\rm a}\mu}(\mathbf r)\,
\delta(\mathbf r-\mathbf R^0_{ia}),\tag{272$'$}\\
\pi_{{\rm a}\mu}(\mathbf r)&=-\frac{i}{\sqrt N}\sum_{\alpha=1}^{3s}\sum_{\mathbf q\in BZ}
\sqrt{\frac{m_{\rm a}\hbar\omega_\alpha(\mathbf q)}{2}}\,e_{{\rm a}\mu\alpha}(\mathbf q)
\big(a_{\mathbf q\alpha}-a^+_{-\mathbf q\alpha}\big)e^{i\mathbf q(\mathbf r-\mathbf r^0_{\rm a})}.\tag{274$'$}
\end{align*}
With (274$'$), $\pi_{{\rm a}\mu}(\mathbf R^0_{ia})$ equals $\pi_{ia\mu}$ of (243). Eqs.~(141), (142),
(237) and (270) follow from the rule and are not displayed here; the final form of (237), which also
carries the factor-$e$ convention of Sec.~\ref{sec:efactor}, is (237$'$) in Sec.~\ref{sec:int-hamiltonian}.
The paramagnetic self-energy (290) becomes
\begin{equation*}
\begin{split}
\Pi^{\rm ion,para}_{\alpha\beta}(\mathbf r,\mathbf r',\omega)={}&-\frac{1}{\hbar c}\frac1N\sum_{\sigma=1}^{3s}
\sum_{\mathbf q\in BZ}\sum_{ia\mu}\sum_{jb\nu}\frac{Z^*_{{\rm a}\alpha\mu}Z^*_{{\rm b}\beta\nu}}{\sqrt{m_{\rm a}m_{\rm b}}}\,
\frac{\hbar\omega^2_\sigma(\mathbf q)}{\omega^2-\omega^2_\sigma(\mathbf q)+i\eta}\\
&\times e_{{\rm a}\mu\sigma}(\mathbf q)e^*_{{\rm b}\nu\sigma}(\mathbf q)\,e^{i\mathbf q(\mathbf R^0_i-\mathbf R^0_j)}\,
\delta(\mathbf r-\mathbf R^0_{ia})\,\delta(\mathbf r'-\mathbf R^0_{jb}).
\end{split}
\tag{290$'$}
\end{equation*}
The lattice sums are now
$\sum_ie^{i\mathbf q\mathbf R^0_i}\delta(\mathbf r-\mathbf R^0_{ia})=\Omega_{\rm uc}^{-1}\sum_{\mathbf G}
e^{i(\mathbf q+\mathbf G)(\mathbf r-\mathbf r^0_{\rm a})}$ and
$\sum_je^{-i\mathbf q\mathbf R^0_j}\delta(\mathbf r'-\mathbf R^0_{jb})=\Omega_{\rm uc}^{-1}\sum_{\mathbf G'}
e^{-i(\mathbf q+\mathbf G')(\mathbf r'-\mathbf r^0_{\rm b})}$, so (293) becomes
\begin{equation*}
\begin{split}
\Pi^{\rm ion,para}_{\alpha\beta}(\mathbf r,\mathbf r',\omega)={}&\sum_{\mathbf G\mathbf G'}\int_{BZ}
\frac{d\mathbf q}{(2\pi)^3}\,e^{i(\mathbf q+\mathbf G)\mathbf r}e^{-i(\mathbf q+\mathbf G')\mathbf r'}
\Big\{-\frac{1}{c\Omega_{\rm uc}}\sum_{\sigma=1}^{3s}\sum_{{\rm a}\mu}\sum_{{\rm b}\nu}\\
&\times\frac{Z^*_{{\rm a}\alpha\mu}Z^*_{{\rm b}\beta\nu}}{\sqrt{m_{\rm a}m_{\rm b}}}\,
\frac{\omega^2_\sigma(\mathbf q)}{\omega^2-\omega^2_\sigma(\mathbf q)+i\eta}\,
e_{{\rm a}\mu\sigma}(\mathbf q)e^{-i(\mathbf q+\mathbf G)\mathbf r^0_{\rm a}}\,
e^*_{{\rm b}\nu\sigma}(\mathbf q)e^{i(\mathbf q+\mathbf G')\mathbf r^0_{\rm b}}\Big\},
\end{split}
\tag{293$'$}
\end{equation*}
and the curly bracket is the $\mathbf G\mathbf G'$ matrix
\begin{align*}
j_{\alpha\sigma}(\mathbf q+\mathbf G)&=\sum_{{\rm a}\mu}\frac{Z^*_{{\rm a}\alpha\mu}}{\sqrt{m_{\rm a}}}\,
e_{{\rm a}\mu\sigma}(\mathbf q)\,e^{-i(\mathbf q+\mathbf G)\mathbf r^0_{\rm a}}
\equiv\sum_{\rm a}j^{\rm a}_{\alpha\sigma}(\mathbf q+\mathbf G),\tag{295$'$}\\
j^*_{\beta\sigma}(\mathbf q+\mathbf G')&=\sum_{{\rm b}\nu}\frac{Z^*_{{\rm b}\beta\nu}}{\sqrt{m_{\rm b}}}\,
e^*_{{\rm b}\nu\sigma}(\mathbf q)\,e^{i(\mathbf q+\mathbf G')\mathbf r^0_{\rm b}},\tag{296$'$}\\
\Pi^{\rm ion,para}_{\alpha\beta,\mathbf G\mathbf G'}(\mathbf q,\omega)&=-\frac{1}{c\Omega_{\rm uc}}
\sum_{\sigma=1}^{3s}\frac{\omega^2_\sigma(\mathbf q)}{\omega^2-\omega^2_\sigma(\mathbf q)+i\eta}\,
j_{\alpha\sigma}(\mathbf q+\mathbf G)\,j^*_{\beta\sigma}(\mathbf q+\mathbf G'),\tag{297$'$}\\
\Pi^{\rm ion,para}_{\alpha\beta,\mathbf G\mathbf G'}(\mathbf q,\omega)&=-\frac{\hbar}{c\Omega_{\rm uc}}
\sum_{\sigma=1}^{3s}\frac{\omega_\sigma(\mathbf q)}{2}\,F_\sigma(\mathbf q,\omega)\,
j_{\alpha\sigma}(\mathbf q+\mathbf G)\,j^*_{\beta\sigma}(\mathbf q+\mathbf G'),\tag{299$'$}\\
\Pi^{\rm ion,dia}_{\alpha\beta}(\mathbf r,\mathbf r',\omega)&=-\frac1c\sum_{ia\mu}\frac{1}{m_{\rm a}}
Z^*_{{\rm a}\alpha\mu}Z^*_{{\rm a}\beta\mu}\,\delta(\mathbf r-\mathbf R^0_{ia})\,\delta(\mathbf r-\mathbf r'),\tag{302$'$}\\
\Pi^{\rm ion,dia}_{\alpha\beta,\mathbf G\mathbf G'}(\mathbf q,\omega)&=-\frac{1}{\Omega_{\rm uc}c}
\sum_{{\rm a}\mu}\frac{1}{m_{\rm a}}Z^*_{{\rm a}\alpha\mu}Z^*_{{\rm a}\beta\mu}\,
e^{-i(\mathbf G-\mathbf G')\mathbf r^0_{\rm a}}.\tag{304$'$}
\end{align*}
Eq.~(306) becomes (297$'$) at $\omega=0$, and then
\begin{align*}
\Pi^{\rm ion,para}_{\alpha\beta,\mathbf G\mathbf G'}(\mathbf q,\omega=0)&=\frac{1}{c\Omega_{\rm uc}}
\sum_{\sigma=1}^{3s}j_{\alpha\sigma}(\mathbf q+\mathbf G)\,j^*_{\beta\sigma}(\mathbf q+\mathbf G'),\tag{307$'$}
\end{align*}
\begin{equation*}
\begin{split}
\Pi^{\rm ion,para}_{\alpha\beta,\mathbf G\mathbf G'}(\mathbf q,\omega=0)={}&\frac{1}{c\Omega_{\rm uc}}
\sum_{{\rm a}\mu}\sum_{{\rm b}\nu}\frac{Z^*_{{\rm a}\alpha\mu}}{\sqrt{m_{\rm a}}}
\frac{Z^*_{{\rm b}\beta\nu}}{\sqrt{m_{\rm b}}}\,e^{-i(\mathbf q+\mathbf G)\mathbf r^0_{\rm a}}\,
e^{i(\mathbf q+\mathbf G')\mathbf r^0_{\rm b}}\\
&\times\sum_{\sigma=1}^{3s}e_{{\rm a}\mu\sigma}(\mathbf q)
e^*_{{\rm b}\nu\sigma}(\mathbf q),
\end{split}
\tag{308$'$}
\end{equation*}
\begin{equation*}
\Pi^{\rm ion,para}_{\alpha\beta,\mathbf G\mathbf G'}(\mathbf q,\omega=0)=\frac{1}{c\Omega_{\rm uc}}
\sum_{{\rm a}\mu}\frac{Z^*_{{\rm a}\alpha\mu}Z^*_{{\rm a}\beta\mu}}{m_{\rm a}}\,
e^{-i(\mathbf G-\mathbf G')\mathbf r^0_{\rm a}}.\tag{309$'$}
\end{equation*}
In (309$'$), completeness (28) keeps only ${\rm a}={\rm b}$, so the $\mathbf q$-dependent phases cancel.
Eqs.~(310) and (311) therefore hold as printed once the subscripts $\mathbf G\mathbf G'$ are added, as
matrix identities. The diagonal $\mathbf G=\mathbf G'$ of (304$'$) is (304) as printed. For a
two-dimensional crystal, (300) and (305) become, with in-plane $\mathbf G$, in-plane positions
$\bm{\rho}^0_{\rm a}$ and heights $z_{\rm a}$,
\begin{equation*}
\begin{split}
\Pi^{\rm ion,para}_{\alpha\beta,\mathbf G\mathbf G'}(\mathbf Q,\omega,z,z')={}&-\frac{\hbar}{cS_{\rm uc}}
\sum_{\sigma=1}^{3s}\frac{\omega_\sigma(\mathbf Q)}{2}F_\sigma(\mathbf Q,\omega)\\
&\times\sum_{\rm ab}j^{\rm a}_{\alpha\sigma}(\mathbf Q+\mathbf G)\,j^{{\rm b}*}_{\beta\sigma}(\mathbf Q+\mathbf G')\,
\delta(z-z_{\rm a})\,\delta(z'-z_{\rm b}),
\end{split}
\tag{300$'$}
\end{equation*}
\begin{equation*}
\begin{split}
\Pi^{\rm ion,dia}_{\alpha\beta,\mathbf G\mathbf G'}(\mathbf Q,\omega,z,z')={}&-\frac{1}{S_{\rm uc}c}
\sum_{{\rm a}\mu}\frac{1}{m_{\rm a}}Z^*_{{\rm a}\alpha\mu}Z^*_{{\rm a}\beta\mu}\\
&\times e^{-i(\mathbf G-\mathbf G')\bm{\rho}^0_{\rm a}}\,\delta(z-z_{\rm a})\,\delta(z'-z_{\rm a}),
\end{split}
\tag{305$'$}
\end{equation*}
where $j^{\rm a}_{\alpha\sigma}(\mathbf Q+\mathbf G)=\sum_\mu Z^*_{{\rm a}\alpha\mu}e_{{\rm a}\mu\sigma}(\mathbf Q)
e^{-i(\mathbf Q+\mathbf G)\bm{\rho}^0_{\rm a}}/\sqrt{m_{\rm a}}$. For a planar sheet ($z_{\rm a}=0$) the atom
sums in (300$'$) collapse to $j_{\alpha\sigma}(\mathbf Q+\mathbf G)j^*_{\beta\sigma}(\mathbf Q+\mathbf G')
\delta(z)\delta(z')$.

\subsection{Corrected and new equations, Secs.~6--8}\label{sec:phase-678}
\begin{align*}
\Pi^{\rm ion}_{jj\mu\nu,\mathbf G\mathbf G'}(\mathbf q,\omega)&=-\frac{\omega^2}{c\Omega_{\rm uc}}
\sum_{\sigma=1}^{3s}\frac{j_{\mu\sigma}(\mathbf q+\mathbf G)\,j^*_{\nu\sigma}(\mathbf q+\mathbf G')}
{\omega^2-\omega^2_\sigma(\mathbf q)+i\eta},\tag{352$'$}\\
\Pi^{\rm ion}_{\rho\rho,\mathbf G\mathbf G'}(\mathbf q,\omega)&=-\frac{c}{\omega^2}\sum_{\mu\nu}
(q_\mu+G_\mu)(q_\nu+G'_\nu)\,\Pi^{\rm ion}_{jj\mu\nu,\mathbf G\mathbf G'}(\mathbf q,\omega).\tag{354a}
\end{align*}
(354a) is inserted after (354); it is the result of pp.~76--77 with the dipoles of (314$'$) and the
substitution (354) modified by the rule above, so (351) still holds. From (380) on, $D$, $\Pi_{jj}$ and
$\Pi^{\rm ion}_{jj}$ are $\mathbf G\mathbf G'$ matrices, as $\Pi^{\rm e}$ already is in (282). The manuscript
keeps only $\mathbf G=\mathbf G'=0$ (p.~95). In that approximation the equations from (385) to (472)
that contain $j_{\mu\sigma}$ only as a symbol keep their printed form, provided that the vertex (the two
unnumbered displays after (385)) is
\begin{equation*}
j_{\mu\sigma}(\mathbf q)=\sum_{{\rm a}\alpha}\frac{Z^*_{{\rm a}\mu\alpha}}{\sqrt{m_{\rm a}}}\,
e_{{\rm a}\alpha\sigma}(\mathbf q)\,e^{-i\mathbf q\mathbf r^0_{\rm a}},\qquad
j^*_{\nu\sigma}(\mathbf q)=\sum_{{\rm b}\beta}\frac{Z^*_{{\rm b}\nu\beta}}{\sqrt{m_{\rm b}}}\,
e^*_{{\rm b}\beta\sigma}(\mathbf q)\,e^{i\mathbf q\mathbf r^0_{\rm b}}.\tag{385$'$}
\end{equation*}
This applies to (389), (391)--(393), (399)--(417), (429), (449)--(454) and, with in-plane
$\bm{\rho}^0_{\rm a}$, (455)--(472). It also applies to the new Eqs.~(410a)--(410h) of
Sec.~\ref{sec:e410} (see Sec.~\ref{sec:int-e410}). With crystal local fields kept, the phonon
self-energy (399) is the $\mathbf G=\mathbf G'=0$ element of the matrix expression (insert after (399),
p.~89)
\begin{equation*}
\begin{aligned}
\Omega^{\rm e}_{\sigma\sigma'}(\mathbf q,\omega)&=-\sum_{\alpha\beta}\sum_{\mathbf G\mathbf G'}
J^*_{\alpha\sigma}(\mathbf q+\mathbf G,\omega)\,D^{\rm e}_{\alpha\beta,\mathbf G\mathbf G'}(\mathbf q,\omega)\,
J_{\beta\sigma'}(\mathbf q+\mathbf G',\omega),\\
J_{\mu\sigma}(\mathbf q+\mathbf G,\omega)&=\sqrt{\frac{\hbar}{2c\Omega_{\rm uc}\omega_\sigma(\mathbf q)}}\,
\omega\,j_{\mu\sigma}(\mathbf q+\mathbf G).
\end{aligned}
\tag{399a}
\end{equation*}
Wherever the vertex is written out, $e^{-i\mathbf q\mathbf r^0_{\rm a}}$ is inserted next to
$e_{{\rm a}\mu\sigma}(\mathbf q)$ and $e^{+i\mathbf q\mathbf r^0_{\rm a}}$ next to $e^*_{{\rm a}\mu\sigma}(\mathbf q)$.
On p.~93 this gives, for the first unnumbered display after (417) [where (295) is inserted], the second
one [where (375) is inserted] and (418),
\begin{equation*}
\begin{split}
g^{\rm dyn,0}_{nm\sigma}(\mathbf k,\mathbf q,\omega)={}&\frac{i}{c\Omega_{\rm uc}}
\frac{\hbar\omega^2}{E_{n\mathbf k}-E_{m\mathbf k+\mathbf q}}\sum_{\alpha\beta=xyz}\sum_{{\rm a}\mu}
\sqrt{\frac{\hbar}{2Nm_{\rm a}\omega_\sigma(\mathbf q)}}\\
&\times Z^*_{{\rm a}\beta\mu}e_{{\rm a}\mu\sigma}(\mathbf q)\,e^{-i\mathbf q\mathbf r^0_{\rm a}}\,
\big[j^\alpha_{n\mathbf k,m\mathbf k+\mathbf q}\big]^*D^0_{\alpha\beta}(\mathbf q,\omega),
\end{split}
\tag{p93a$'$}
\end{equation*}
\begin{equation*}
\begin{split}
g^{\rm dyn,0}_{nm\sigma}(\mathbf k,\mathbf q,\omega)={}&\frac{4\pi i}{\Omega_{\rm uc}}
\frac{\hbar}{E_{n\mathbf k}-E_{m\mathbf k+\mathbf q}}\sum_{\alpha\beta=xyz}\sum_{{\rm a}\mu}
\sqrt{\frac{\hbar}{2Nm_{\rm a}\omega_\sigma(\mathbf q)}}\\
&\times Z^*_{{\rm a}\beta\mu}e_{{\rm a}\mu\sigma}(\mathbf q)\,e^{-i\mathbf q\mathbf r^0_{\rm a}}\,
\frac{-q_\alpha q_\beta+\frac{\omega^2}{c^2}\delta_{\alpha\beta}}{q^2-\frac{\omega^2}{c^2}-i\eta}
\big[j^\alpha_{n\mathbf k,m\mathbf k+\mathbf q}\big]^*,
\end{split}
\tag{p93b$'$}
\end{equation*}
\begin{equation*}
\begin{split}
\lim_{c\to\infty}g^{\rm dyn,0}_{nm\sigma}(\mathbf k,\mathbf q,\omega)={}&\frac{4\pi i}{\Omega_{\rm uc}}
\sum_{\beta=xyz}\sum_{{\rm a}\mu}\sqrt{\frac{\hbar}{2Nm_{\rm a}\omega_\sigma(\mathbf q)}}\,
Z^*_{{\rm a}\beta\mu}e_{{\rm a}\mu\sigma}(\mathbf q)\,e^{-i\mathbf q\mathbf r^0_{\rm a}}\,\frac{q_\beta}{q^2}\\
&\times\sum_{\alpha=xyz}\frac{-\hbar q_\alpha}{E_{n\mathbf k}-E_{m\mathbf k+\mathbf q}}
\big[j^\alpha_{n\mathbf k,m\mathbf k+\mathbf q}\big]^*,
\end{split}
\tag{418$'$}
\end{equation*}
and on pp.~96 and 98
\begin{equation*}
\begin{split}
g^{\rm dyn}_{nm\sigma}(\mathbf k,\mathbf q,\omega)={}&\frac{i}{c\Omega_{\rm uc}}
\frac{\hbar\omega^2}{E_{n\mathbf k}-E_{m\mathbf k+\mathbf q}}\sum_{{\rm a}\mu}\sum_{\alpha\beta=xyz}
\sqrt{\frac{\hbar}{2Nm_{\rm a}\omega_\sigma(\mathbf q)}}\\
&\times Z^*_{{\rm a}\alpha\mu}e^*_{{\rm a}\mu\sigma}(\mathbf q)\,e^{i\mathbf q\mathbf r^0_{\rm a}}\,
D^{\rm e}_{\alpha\beta}(\mathbf q,\omega)\,j^\beta_{n\mathbf k,m\mathbf k+\mathbf q},
\end{split}
\tag{430$'$}
\end{equation*}
\begin{equation*}
\begin{split}
\lim_{c\to\infty}g^{\rm dyn}_{nm\sigma}(\mathbf k,\mathbf q,\omega)={}&-\frac{4\pi i}{\Omega_{\rm uc}}
\frac{\hbar}{E_{n\mathbf k}-E_{m\mathbf k+\mathbf q}}\sum_{{\rm a}\mu}\sum_{\alpha\beta=xyz}
\sqrt{\frac{\hbar}{2Nm_{\rm a}\omega_\sigma(\mathbf q)}}\\
&\times Z^*_{{\rm a}\alpha\mu}e^*_{{\rm a}\mu\sigma}(\mathbf q)\,e^{i\mathbf q\mathbf r^0_{\rm a}}\,
\frac{q_\alpha q_\beta}{\mathbf q\cdot\hat\epsilon^{\rm e}(\mathbf q,\omega)\cdot\mathbf q}\,
j^\beta_{n\mathbf k,m\mathbf k+\mathbf q},
\end{split}
\tag{447$'$}
\end{equation*}
\begin{equation*}
\begin{split}
\lim_{c\to\infty}g^{\rm dyn}_{nm\sigma}(\mathbf k,\mathbf q,\omega)={}&\frac{4\pi i}{\Omega_{\rm uc}}
\sum_{{\rm a}\mu}\sum_\alpha\sqrt{\frac{\hbar}{2Nm_{\rm a}\omega_\sigma(\mathbf q)}}\,
Z^*_{{\rm a}\alpha\mu}e^*_{{\rm a}\mu\sigma}(\mathbf q)\,e^{i\mathbf q\mathbf r^0_{\rm a}}\\
&\times\frac{q_\alpha}{\mathbf q\cdot\hat\epsilon^{\rm e}(\mathbf q,\omega)\cdot\mathbf q}\sum_\beta\frac{-\hbar q_\beta j^\beta_{n\mathbf k,m\mathbf k+\mathbf q}}{E_{n\mathbf k}-E_{m\mathbf k+\mathbf q}}.
\end{split}
\tag{448$'$}
\end{equation*}
These displays contain the electronic current $j^\mu_{n\mathbf k,m\mathbf k+\mathbf q}$, which carries the
charge, and are not affected by the factor-$e$ correction. The displays on pp.~95, 96 and 98, obtained
from (418$'$) and (448$'$) with the charge--current relation, are affected by both corrections. Their
final forms are (p95$'$), (p96$'$) and (p98$'$) in Sec.~\ref{sec:int-vertices}.

Eq.~(454) keeps its form with (385$'$). For an insulator with
$W^{\rm e}_{\mathbf G\mathbf G'}=4\pi\delta_{\mathbf G\mathbf G'}/[(\mathbf q+\mathbf G)\cdot\hat\epsilon^\infty\cdot
(\mathbf q+\mathbf G)]$, (399a) gives the matrix, $\mathbf G$-resolved form of (454) (placement on p.~99:
Sec.~\ref{sec:int-p99})
\begin{equation*}
\begin{split}
\tilde\Omega^2_{{\rm a}\mu,{\rm b}\nu}(\mathbf q)={}&\tilde\Omega^2_{0,{\rm a}\mu,{\rm b}\nu}(\mathbf q)\\
&+\frac{4\pi}{\Omega_{\rm uc}\sqrt{m_{\rm a}m_{\rm b}}}\sum_{\mathbf G}
\frac{[(\mathbf q+\mathbf G)\hat Z^*_{\rm a}]_\mu\,[(\mathbf q+\mathbf G)\hat Z^*_{\rm b}]_\nu}
{(\mathbf q+\mathbf G)\cdot\hat\epsilon^\infty\cdot(\mathbf q+\mathbf G)}\,
e^{i(\mathbf q+\mathbf G)(\mathbf r^0_{\rm a}-\mathbf r^0_{\rm b})},
\end{split}
\tag{454a}
\end{equation*}
with $[(\mathbf q+\mathbf G)\hat Z^*_{\rm a}]_\mu=\sum_\alpha(q_\alpha+G_\alpha)Z^*_{{\rm a}\alpha\mu}$ and
$\tilde\Omega^2_0(\mathbf q)$ the bare (short-range) dynamical matrix, whose eigenvectors are
$e_{{\rm a}\mu\sigma}(\mathbf q)$ and eigenvalues $\omega^2_\sigma(\mathbf q)$. Projected on mode $\sigma$ and
restricted to $\mathbf G=0$, (454a) gives (454). The full sum is the dipole--dipole matrix of X.~Gonze
and C.~Lee, Phys.\ Rev.\ B \textbf{55}, 10355 (1997), once the on-site self term is removed by Ewald
regularization. With every dipole at the cell origin, the ${\rm a}\neq{\rm b}$ terms would not converge
either.

\subsection{What is unchanged, and what changes}
At $\mathbf G=0$ only the relative phases $e^{-i\mathbf q(\mathbf r^0_{\rm a}-\mathbf r^0_{\rm b})}
=1+O(q|\mathbf r^0_{\rm a}-\mathbf r^0_{\rm b}|)$ are observable: a common factor $e^{-i\mathbf q\mathbf c}$ (a
shift $\mathbf c$ of the origin) drops out of $|g|^2$ and of $j_{\alpha\sigma}j^*_{\beta\sigma}$. The following
are therefore unchanged:
\begin{itemize}
\item the $\mathbf q\to0$, $\mathbf G=0$ limit of (131$'$) for optical modes, i.e.\ Fr\"ohlich's vertex
(Giustino's Eq.~(7));
\item Sec.~4.2.1, Eqs.~(144)--(165), where the field is uniform over a cell;
\item (406)--(410) at $\mathbf q=0$, and (410$'$), (410e) of Sec.~\ref{sec:e410};
\item the LO--TO splitting, i.e.\ the $\mathbf q\to0$ limit of (454);
\item the polariton region, where $q\lesssim\omega_{\rm TO}\sqrt{\epsilon(0)}/c\approx6\times10^{-5}$~\AA$^{-1}$
in GaAs [$\epsilon(0)$ the static dielectric constant, as in (165)], $q|\mathbf r^0_{\rm b}-\mathbf r^0_{\rm a}|
\approx10^{-4}$, and the oscillator strength changes by $5\times10^{-9}$;
\item the sum rule (309) and Eqs.~(310)--(311) on the diagonal $\mathbf G=\mathbf G'$;
\item the form of (175)--(180) and of (235) [up to the factor $e$ of Sec.~\ref{sec:efactor}], and all
electronic and photon-propagator equations.
\end{itemize}
The correction matters in three places:
\begin{itemize}
\item at finite $q$ in cells with $s\ge2$. In GaAs ($|\mathbf r^0_{\rm As}-\mathbf r^0_{\rm Ga}|=2.45$~\AA) the
phase difference $\mathbf q(\mathbf r^0_{\rm As}-\mathbf r^0_{\rm Ga})$ is 0.24~rad at $q=0.1$~\AA$^{-1}$ and
1.22~rad at 0.5~\AA$^{-1}$ along [111]. In the central-force model, with eigenvectors computed at finite
$q$, $|g^{\mathcal L}_{\rm LO}|^2$ and the dipolar shift in (454) change by 0.5--0.8\% at
$q=0.1$~\AA$^{-1}$ and by 12--22\% at $q=0.5$~\AA$^{-1}$, up or down depending on the direction and on
the cell choice of the printed form;
\item in every $\mathbf G\neq0$ element, where $e^{-i\mathbf G\mathbf r^0_{\rm a}}$ survives even at
$\mathbf q=0$;
\item for acoustic modes, see (131b), and at finite $q$ in (410c) of Sec.~\ref{sec:e410}.
\end{itemize}

\subsection{Replacement prose}
\begin{itemize}
\item p.~24, before (121): as given at the start of the corrected equations of Sec.~3 above.
\item p.~25, after (131), replace ``which is exactly equal to Eq.4 \ldots\ dimension of charge.'' by:
\begin{quote}
This is Eq.~(180) of F.~Giustino, Rev.\ Mod.\ Phys.\ \textbf{89}, 015003 (2017), the corrected form of
Eq.~(4) of C.~Verdi and F.~Giustino, Phys.\ Rev.\ Lett.\ \textbf{115}, 176401 (2015). The basis phase
$e^{-i(\mathbf q+\mathbf G)\mathbf r^0_{\rm a}}$ is required because the phonon eigenvectors carry the cell
phase $e^{i\mathbf q\mathbf R^0_i}$ of Eq.~(83). Verdi and Giustino's Eq.~(4) lacked it (Giustino's
footnote~7, which notes that it is needed for the acoustic modes near $\mathbf q=0$). The remaining
differences are the factor $e$ hidden in $\hat Z^*$, which has the dimension of charge, the band and
momentum labels, and the exclusion of $\mathbf G=-\mathbf q$.
\end{quote}
\item p.~62, after (290), replace ``The above form \ldots\ avoided whenever possible.'' by:
\begin{quote}
The above form of the current propagator follows from the fact that each atomic site $\mathbf R^0_{ia}$
carries a polarizable point dipole. The periodic arrangement of these dipoles produces Bragg components
$\mathbf q+\mathbf G$, weighted by the structure factors $e^{-i(\mathbf q+\mathbf G)\mathbf r^0_{\rm a}}$ of the
vertices (295$'$). They are reliable only for $|\mathbf q+\mathbf G|$ small compared with the inverse size of
the charge induced around an atom, because an atom is not a point dipole: quadrupoles enter at the next
order. In practice the $\mathbf G$ sum is cut off smoothly, e.g.\ with $e^{-a|\mathbf q+\mathbf G|^2}$ as in
Verdi and Giustino.
\end{quote}
\item p.~63, replace ``We see that ionic self-energy is a constant in $\mathbf G\mathbf G'$ space \ldots\
dipole grating.'' by:
\begin{quote}
The ionic self-energy depends on $\mathbf G$ and $\mathbf G'$ only through the basis phases of the ionic
current vertices. It is the response of $s$ interpenetrating gratings of point dipoles, one per
sublattice, each displaced by $\mathbf r^0_{\rm a}$ from the Bravais lattice. Its $\mathbf G=\mathbf G'=0$
element is the macroscopic response used in Secs.~6.1--8.
\end{quote}
\item p.~64, after (299$'$): ``where the $\mathbf G\mathbf G'$ dependence sits in the vertices.'' After
(300$'$): ``Within the point-dipole model this is the exact polarizability of a planar sheet.''
\item p.~85, after (380), add: ``Here and below only the $\mathbf G=\mathbf G'=0$ element is kept. This
neglects the crystal local fields of the electrons and, through the structure factors of (295$'$), those
of the ions.''
\item p.~86, after (385): the ionic current vertices are (385$'$). p.~89, after (399): ``With crystal local
fields kept, (399) is the $\mathbf G=\mathbf G'=0$ element of'' (399a). p.~93, after (417): the two
displays are (p93a$'$) and (p93b$'$), and (418) becomes (418$'$).
\item pp.~95--99: Sec.~\ref{sec:int-vertices}; p.~99 after (454): Sec.~\ref{sec:int-p99}.
\end{itemize}

\subsection{Justification}
The normal-mode expansion (7), (83) gives every atom of cell $i$ the Bloch phase $e^{i\mathbf q\mathbf R^0_i}$
of its cell. The Fourier component at $\mathbf q+\mathbf G$ of the dipole density of mode $(\mathbf q,\alpha)$
is then $\Omega_{\rm uc}^{-1}\sum_{\rm a}\hat Z^*_{\rm a}\mathbf e_{{\rm a}\alpha}(\mathbf q)
e^{-i(\mathbf q+\mathbf G)\mathbf r^0_{\rm a}}/\sqrt{m_{\rm a}}$ times the phonon amplitude. Sec.~2 already
obtains exactly this structure for the pseudopotential, Eq.~(101). Dropping the factor has three
consequences:
\begin{itemize}
\item the vertex depends on the arbitrary assignment of atoms to cells, Eq.~(131a);
\item it predicts a spurious, cell-dependent long-range coupling to acoustic phonons, $|g|^2\propto1/q$,
even in centrosymmetric crystals;
\item in the ionic self-energy and in the dipole--dipole dynamical matrix (454a), the $\mathbf G$ sums of
the printed form describe coincident dipoles. The on-site terms ${\rm a}={\rm b}$ need regularization
(removal of the Ewald self term) in both forms; with coincident dipoles the ${\rm a}\neq{\rm b}$ terms
diverge as well. The vertex $g^{\mathcal L}$ itself stays finite, because coincident dipoles simply add
into one dipole per cell.
\end{itemize}
With the factor, (131$'$) coincides with Giustino's Eq.~(180), and the ionic self-energy
(297$'$)--(304$'$) reproduces the Ewald-summed dipole lattice of Gonze and Lee, Eq.~(454a). This was
checked numerically: the $\mathbf G$ sum is independent of the Ewald parameter only with the sign
$e^{-i(\mathbf q+\mathbf G)\mathbf r^0_{\rm a}}$. The cell-choice invariance (131a), the acoustic expansion
(131b) (remainder of order $q^2$), the matrix sum rule (310) with (304$'$) and (309$'$), and the GaAs
numbers above were also checked numerically. All $\mathbf q\to0$, $\mathbf G=0$ results of the manuscript for
optical modes are untouched.

%======================================================================================
\section{E.410 is the $\mathbf q\to0$ upper phonon polariton, not the LO--TO splitting (Sec.~7.1)}\label{sec:e410}

\subsection{Location}
Sec.~7.1, pp.~90--91: the title, the first sentences, the text around Eqs.~(406)--(410), and Eq.~(410);
the title of Sec.~7.3 and the text after Eq.~(454) on p.~99.

\subsection{Problem}
The title of Sec.~7.1 and its first sentence present Eq.~(410) as the LO--TO splitting, ``completely
analogous to the non-analytic dynamical matrix correction done by Quantum Espresso''. The algebra of
Eqs.~(406)--(410) is correct, but Eq.~(410) is a different limit. It is taken at $\mathbf q=0$ for finite
$c$, where the vacuum propagator, Eq.~(408), is isotropic. Every polar mode, transverse and longitudinal,
is therefore shifted by the same $\frac{4\pi}{\Omega_{\rm uc}}\sum_\alpha|j_{\alpha\sigma}|^2$, and there is no
LO--TO splitting at this point. Eq.~(410) is the $\mathbf q\to0$ end point of the upper phonon polariton
(Huang) of a crystal without electrons. The standard non-analytic term that Quantum Espresso adds to
the dynamical matrix is $\propto|\hat{\mathbf q}\cdot\mathbf j_\sigma|^2/(\hat{\mathbf q}\cdot\hat\epsilon^\infty
\cdot\hat{\mathbf q})$: it leaves the TO modes unshifted and carries $1/\epsilon_\infty$. Its counterpart in
the manuscript is Eq.~(454), obtained with $c\to\infty$ taken first and the electrons resummed into
$D^{\rm e}$. Eq.~(410) differs from (454) in two independent ways, which must be kept apart:
\begin{itemize}
\item[(a)] \emph{Order of limits.} For the longitudinal mode, Eq.~(410) is the unscreened LO frequency
whichever limit is taken first, and coincides with Eq.~(454) for $\epsilon_\infty=1$. The order of limits
changes only the transverse modes, which at $\mathbf q\to0$ and finite $c$ are degenerate with the
longitudinal one. For GaAs only the TO modes move: 35.9~meV $\to$ 33.3~meV.
\item[(b)] \emph{Electronic screening.} For GaAs, Eq.~(410) gives 55.6~meV, against the measured LO
energy of 36.2~meV. This gap comes from the electron-free approximation, which Sec.~7.1 states
(``Suppose that we neglect the interaction with electrons''), not from the order of limits. With
electrons, the same $\mathbf q\to0$, finite-$c$ limit puts all three polar modes at 35.9~meV
[Eq.~(410e) below]: 55.6~meV $\to$ 35.9~meV.
\end{itemize}

\subsection{What changes, what does not}
No displayed equation is algebraically wrong, and no later equation uses Eq.~(410), so no downstream
equation inherits a change. \emph{Changed:} the title of Sec.~7.1, its opening sentences, the sentences
around Eqs.~(406)--(410) (including two photon/phonon slips and the subscript slip
$\Omega^0_\sigma\to\Omega^0_{\sigma\sigma}$ in (406)), the left-hand side of Eq.~(410) [now (410$'$)], a new
paragraph and a new Sec.~7.1.1 with Eqs.~(410a)--(410h) after Eq.~(410$'$), and the title of Sec.~7.3
together with one paragraph after Eq.~(454). \emph{Unchanged:} Eqs.~(406)--(409); Eqs.~(411)--(454),
apart from the basis phase of Sec.~\ref{sec:phase}; Eqs.~(455)--(459) and the $z$-dependent Dyson
equation on p.~100, which are the general, finite-$c$ transcription to 2D; the electrostatic
Eqs.~(460)--(472) of Sec.~8.1, which do not describe $Q\lesssim\omega/c$ but need no change; and the
classical discussion of pp.~33--35, Eqs.~(159)--(165), which assumes $\nabla\times\mathbf E=0$ and is the
electrostatic picture.

\subsection{Corrected and new equations, replacement prose}
\paragraph{1. Title of Sec.~7.1 (p.~90).}
Replace ``Lowest self-energy correction and LO-TO splitting'' by ``Lowest self-energy correction and the
phonon polariton at $\mathbf q=0$''.

\paragraph{2. Opening of Sec.~7.1 (p.~90).}
Replace ``This procedure is completely analogous to the non-analytic dynamical matrix correction done by
Quantum Espresso. Let's continue with the diagonal approximation. Suppose that we neglect the interaction
with electrons, then'' by:
\begin{quote}
The self-energy (399) describes how the long-range electromagnetic field of the Born-charge dipoles of a
polar mode acts back on the bare, short-range phonons. In this respect it plays the role of the standard
non-analytic correction that Quantum Espresso adds to the DFPT dynamical matrix [\textit{P.~Giannozzi,
S.~de~Gironcoli, P.~Pavone and S.~Baroni, Phys.\ Rev.\ B} \textbf{43}, 7231 (1991); \textit{S.~Baroni,
S.~de~Gironcoli, A.~Dal~Corso and P.~Giannozzi, Rev.\ Mod.\ Phys.} \textbf{73}, 515 (2001); \textit{X.~Gonze
and C.~Lee, Phys.\ Rev.\ B} \textbf{55}, 10355 (1997)]. Here, however, the field is retarded, and the
result depends on the order in which the limits $\mathbf q\to0$ and $c\to\infty$ are taken. We first
evaluate the lowest, electron-free correction at $\mathbf q=0$ for finite $c$. General $\mathbf q$ and
electronic screening follow in Sec.~7.1.1, and the electrostatic limit, which gives the non-analytic
correction, in Sec.~7.3. We keep the diagonal approximation and neglect the interaction with electrons,
so that
\end{quote}
In the next lines, ``the photon self-energy is'' (before (406)) should read ``the phonon self-energy is'',
with $\Omega^0_{\sigma\sigma}$ instead of $\Omega^0_\sigma$ on the left of Eq.~(406); and ``the new corrected
photon energy satisfy'' (before (407)) should read ``the new corrected phonon energy satisfies''.

\paragraph{3. Around Eqs.~(408)--(410) (p.~91).}
Replace ``After the photon propagator (375) in optical limit $\mathbf q=0$'' by ``After the photon
propagator (375), evaluated at $\mathbf q=0$ for finite $c$, i.e.\ deep inside the light cone
($cq\ll\omega$), where it is isotropic,''. Replace ``From this is then possible to explicitly derive the
phonon energy correction in the optical limit'' by ``The pole condition (407) then gives, for
$\mathbf q\to0$ at finite $c$ and without electrons,'' and write Eq.~(410) as
\begin{equation*}
\tilde\omega^2_\sigma(\mathbf q\to0)\Big|_{c\ \text{finite}}
=\omega^2_\sigma(\mathbf q=0)+\frac{4\pi}{\Omega_{\rm uc}}\sum_\alpha\big|j_{\alpha\sigma}(\mathbf q=0)\big|^2 .
\tag{410$'$}
\end{equation*}

\paragraph{4. Insert after Eq.~(410$'$) (p.~91).}
\begin{quote}
Equation (410$'$) is not the LO--TO splitting. Because $D^0(\mathbf q=0,\omega)$ is isotropic, transverse
and longitudinal polar modes are shifted by the same amount, so there is no splitting at this point, and
since electrons are absent the shift is unscreened ($\epsilon_\infty=1$). Equation (410$'$) is the common
$\mathbf q\to0$ frequency of the longitudinal mode and of the upper transverse phonon polariton
[\textit{K.~Huang, Proc.\ R.\ Soc.\ Lond.\ A} \textbf{208}, 352 (1951)] of a crystal without electrons; the
lower polariton branch is photon-like and tends to $\omega=0$. The LO--TO splitting given by the standard
non-analytic correction, a term $\propto|\hat{\mathbf q}\cdot\mathbf j_\sigma|^2/(\hat{\mathbf q}\cdot
\hat\epsilon^\infty\cdot\hat{\mathbf q})$, leaves the TO modes unshifted and carries $1/\epsilon_\infty$. Its
counterpart here is Eq.~(454), obtained with $c\to\infty$ taken first and the electrons resummed into
$D^{\rm e}$. For the longitudinal mode alone the two agree when $\epsilon_\infty=1$: Eq.~(410$'$) is then
the electron-free LO frequency, whichever limit is taken first. Both limits, and the separate role of
electronic screening, follow from one formula, derived next.
\end{quote}

\noindent New Sec.~7.1.1, \emph{Finite $\mathbf q$, electronic screening and the order of limits}:
\begin{quote}
Keep $\mathbf q$ finite, write $\hat{\mathbf q}=\mathbf q/q$, and model the electrons of an insulator by the
low-frequency, $\mathbf G=0$ form of the correlator (383),
$\Pi^{\rm e,0}_{jj\mu\nu}(\mathbf q,\omega)\simeq\frac{(\epsilon_\infty-1)\,\omega^2}{4\pi c}\,\delta_{\mu\nu}$, so
that $\hat{\mathbf q}\cdot\hat\epsilon^{\rm e}\cdot\hat{\mathbf q}=\epsilon_\infty$ in Eq.~(432). The longitudinal
and transverse projections of Eq.~(375) are $-4\pi c/\omega^2$ and $-4\pi c/(\omega^2-c^2q^2)$, the two
blocks decouple, and the Dyson equation (388) gives
\begin{equation*}
D^{\rm e}_{\mu\nu}(\mathbf q,\omega)=-4\pi c\left[\frac{\hat q_\mu\hat q_\nu}{\epsilon_\infty\,\omega^2}
+\frac{\delta_{\mu\nu}-\hat q_\mu\hat q_\nu}{\epsilon_\infty\,\omega^2-c^2q^2}\right].
\tag{410a}
\end{equation*}
For $\epsilon_\infty=1$ this is Eq.~(375). For an anisotropic tensor $\hat\epsilon^\infty$,
$D^{\rm e}_{\mu\nu}=-4\pi c\,(\Lambda^{-1})_{\mu\nu}$ with the Maxwell matrix
$\Lambda_{\mu\nu}(\mathbf q,\omega)=\omega^2\epsilon^\infty_{\mu\nu}-c^2(q^2\delta_{\mu\nu}-q_\mu q_\nu)$.
Inserting (410a) and the vertices (392), with $j_{\alpha\sigma}(\mathbf q)$ the ionic vertex (385$'$), which
carries the basis phase $e^{-i\mathbf q\mathbf r^0_{\rm a}}$, in Eq.~(399) gives
\begin{equation*}
\Omega^{\rm e}_{\sigma\sigma}(\mathbf q,\omega)=\frac{2\pi\hbar}{\Omega_{\rm uc}\,\omega_\sigma(\mathbf q)}
\left[\frac{|\hat{\mathbf q}\cdot\mathbf j_\sigma|^2}{\epsilon_\infty}
+\frac{\omega^2\,|(1-\hat{\mathbf q}\hat{\mathbf q})\cdot\mathbf j_\sigma|^2}{\epsilon_\infty\,\omega^2-c^2q^2}\right],
\tag{410b}
\end{equation*}
with $|\hat{\mathbf q}\cdot\mathbf j_\sigma|^2=\big|\sum_\alpha\hat q_\alpha j_{\alpha\sigma}(\mathbf q)\big|^2$ and
$|(1-\hat{\mathbf q}\hat{\mathbf q})\cdot\mathbf j_\sigma|^2=\sum_\alpha|j_{\alpha\sigma}(\mathbf q)|^2
-|\hat{\mathbf q}\cdot\mathbf j_\sigma|^2$. For $\epsilon_\infty=1$, (410b) is $\Omega^0_{\sigma\sigma}$ of
Eq.~(406); the difference is $\Delta\Omega^{\rm e}_{\sigma\sigma}$ of Eq.~(402), which at $\mathbf q=0$ equals
$\frac{2\pi\hbar}{\Omega_{\rm uc}\omega_\sigma}(\epsilon_\infty^{-1}-1)\sum_\alpha|j_{\alpha\sigma}|^2<0$. The poles
of Eq.~(405) satisfy
\begin{equation*}
\omega^2=\omega^2_\sigma(\mathbf q)+\frac{4\pi}{\Omega_{\rm uc}}
\left[\frac{|\hat{\mathbf q}\cdot\mathbf j_\sigma|^2}{\epsilon_\infty}
+\frac{\omega^2\,|(1-\hat{\mathbf q}\hat{\mathbf q})\cdot\mathbf j_\sigma|^2}{\epsilon_\infty\,\omega^2-c^2q^2}\right].
\tag{410c}
\end{equation*}
Without the diagonal approximation, the $3s\times3s$ Dyson equation (398) gives instead
\begin{equation*}
\det\left[(\omega^2-\omega^2_\sigma)\,\delta_{\sigma\sigma'}
-\frac{4\pi\omega^2}{\Omega_{\rm uc}}\sum_{\alpha\beta}j^*_{\alpha\sigma}(\mathbf q)\,(\Lambda^{-1})_{\alpha\beta}\,
j_{\beta\sigma'}(\mathbf q)\right]=0 .
\tag{410d}
\end{equation*}
For a single polar mode, or for degenerate modes chosen with mutually orthogonal $\mathbf j_\sigma$ along
or across $\hat{\mathbf q}$ (isotropic $\epsilon_\infty$), (410d) reduces to (410c). The longitudinal term
of (410c) contains no $c$: the LO frequency is not retarded. The transverse term contains $c$ and $q$
only through $cq$, so the limits $\mathbf q\to0$ and $c\to\infty$ do not commute. For $\mathbf q\to0$ at
finite $c$,
\begin{equation*}
\tilde\omega^2_\sigma=\omega^2_\sigma+\frac{4\pi}{\Omega_{\rm uc}\,\epsilon_\infty}\sum_\alpha|j_{\alpha\sigma}|^2 ,
\tag{410e}
\end{equation*}
while for $cq\gg\sqrt{\epsilon_\infty}\,\omega$
\begin{equation*}
\begin{aligned}
\tilde\omega^2_\sigma(\mathbf q)&=\omega^2_\sigma(\mathbf q)+\frac{4\pi}{\Omega_{\rm uc}\,\epsilon_\infty}\,
|\hat{\mathbf q}\cdot\mathbf j_\sigma|^2\\
&=\omega^2_\sigma(\mathbf q)+\frac{W^{\rm e}(\mathbf q,0)}{\Omega_{\rm uc}}\Big|\sum_\alpha q_\alpha
j_{\alpha\sigma}(\mathbf q)\Big|^2 .
\end{aligned}
\tag{410f}
\end{equation*}
Eq.~(410e) is isotropic and reduces to Eq.~(410$'$) for $\epsilon_\infty=1$. Eq.~(410f), with
$W^{\rm e}(\mathbf q,0)=4\pi/(q^2\epsilon_\infty)$, is Eq.~(454); with a tensor $\hat\epsilon^\infty$ it becomes
$4\pi|\hat{\mathbf q}\cdot\mathbf j_\sigma|^2/(\Omega_{\rm uc}\,\hat{\mathbf q}\cdot\hat\epsilon^\infty\cdot
\hat{\mathbf q})$, the standard non-analytic term, and in that limit (410d) is the mode-basis form of the
dynamical matrix plus the non-analytic correction. For a longitudinal mode the two limits give the same
frequency; they differ only for the transverse modes. For a cubic diatomic crystal
($Z^*_{\rm a}=-Z^*_{\rm b}=Z^*$, $\sum_\alpha|j_{\alpha\sigma}|^2=Z^{*2}/M$ for each optical mode, $M$ the
reduced mass, $\omega_p^2=4\pi Z^{*2}/(\Omega_{\rm uc}M)$ the ionic plasma frequency defined after
Eq.~(152), p.~33), (410c) gives one longitudinal mode and two transverse polariton branches,
\begin{equation*}
\begin{gathered}
\omega^2_{\rm L}=\omega^2_\sigma+\frac{\omega_p^2}{\epsilon_\infty},\qquad
c^2q^2=\omega^2\epsilon(\omega),\quad \epsilon(\omega)=\epsilon_\infty+\frac{\omega_p^2}{\omega^2_\sigma-\omega^2},\\
\omega^2_\pm=\frac{B\pm\sqrt{B^2-4\epsilon_\infty c^2q^2\omega^2_\sigma}}{2\epsilon_\infty},\qquad
B=\epsilon_\infty\omega^2_{\rm L}+c^2q^2 .
\end{gathered}
\tag{410g}
\end{equation*}
At $\mathbf q\to0$, $\omega_+\to\omega_{\rm L}$ and $\omega_-\to0$; for $cq\gg\sqrt{\epsilon_\infty}\,\omega_\sigma$,
$\omega_-\to\omega_\sigma=\omega_{\rm TO}$ and $\omega_{\rm L}=\omega_{\rm LO}$. The crossover between the two
regimes is at
\begin{equation*}
q_c=\frac{\sqrt{\epsilon_\infty}\,\omega}{c},\qquad \omega\simeq\omega_\sigma .
\tag{410h}
\end{equation*}
For GaAs ($Z^*=2.07\,|e|$, $\Omega_{\rm uc}=45.2$~\AA$^3$, $M=36.1$~amu, $\epsilon_\infty=10.9$,
$\hbar\omega_{\rm TO}=33.3$~meV, $\hbar\omega_{\rm LO}=36.2$~meV; literature values) one finds
$\hbar\omega_p=44.6$~meV. With $\omega_\sigma(0)=\omega_{\rm TO}$, since the $\mathbf G=0$ treatment has no
Lorentz local-field term, Eq.~(410$'$) gives 55.6~meV, (410e) gives 35.9~meV for all three polar modes,
and (410f) gives LO at 35.9~meV and TO at 33.3~meV. The difference between 55.6 and 35.9~meV is
electronic screening; the order of limits decides only whether the transverse modes share the shift
(35.9~meV) or stay at 33.3~meV. The crossover is $q_c=5.6\times10^{-5}$~\AA$^{-1}$, about
$5\times10^{-5}$ of $2\pi/a$. At $q=10\,q_c$ the TO mode lies only 0.03~meV below $\omega_{\rm TO}$. Except
for $q\lesssim q_c$, retardation is therefore negligible, and the electrostatic result (454) applies within
the manuscript's $\mathbf G=0$ treatment (in the diagonal approximation; for mixed polar modes, the
$c\to\infty$ limit of (410d)). At finite $\mathbf q$, $j_{\alpha\sigma}(\mathbf q)$ is the vertex (385$'$) with
the basis phase $e^{-i\mathbf q\mathbf r^0_{\rm a}}$; without it, (410c) would give the acoustic modes a
spurious, cell-dependent dipolar shift [cf.\ (131b)]. Eq.~(410$'$) and its screened form (410e) describe
the polariton region $q\lesssim q_c$ ($q\lesssim\omega/c$ without electrons), probed by infrared
spectroscopy and near-forward (polariton) Raman scattering. Visible-light backscattering Raman probes
$q\approx10^{-2}$~\AA$^{-1}$, more than $10^2\,q_c$, and sees the split TO and LO of (410f). In the
polariton region the direction dependence of the electrostatic term at $\Gamma$ is lifted by
retardation.
\end{quote}

\paragraph{5. Sec.~7.3 (p.~99).}
Title: ``Phonon's self-energy in electrostatic limit and the LO--TO splitting''. After the sentence that
follows Eq.~(454), i.e.\ after ``\ldots\ Nano Lett.\ \textbf{17}, 3758 (2017).'', insert (first of the two
insertions after (454); order in Sec.~\ref{sec:int-p99}):
\begin{quote}
With $W^{\rm e}(\mathbf q,0)=4\pi/(\mathbf q\cdot\hat\epsilon^\infty\cdot\mathbf q)$, Eq.~(454) is the standard
non-analytic term that Quantum Espresso adds to the DFPT dynamical matrix (references in Sec.~7.1). It
shifts only modes with $\hat{\mathbf q}\cdot\mathbf j_\sigma\neq0$, and the shift is reduced by the factor
$1/(\hat{\mathbf q}\cdot\hat\epsilon^\infty\cdot\hat{\mathbf q})$ relative to the unscreened one. It is the
$cq\gg\sqrt{\epsilon_\infty}\,\omega$ limit (410f) of Eq.~(410c), whereas Eq.~(410$'$) is its electron-free
$\mathbf q\to0$ limit at finite $c$.
\end{quote}

\subsection{Justification}
Equations (410a)--(410d) use only the manuscript's definitions. (410a) solves Eq.~(388) with Eq.~(375).
(410b) is Eq.~(399) with the vertices (392). (410c) is the pole condition of Eq.~(405), and (410d)
follows from Eq.~(398). The checks are as follows.
(i)~Dimensions: in Gaussian units $[Z^{*2}]$ is energy$\times$length, so
$\frac{4\pi}{\Omega_{\rm uc}}|j_{\alpha\sigma}|^2$ is a squared frequency, $\Omega^{\rm e}$ an energy and $q_c$ an
inverse length.
(ii)~$\epsilon_\infty=1$, $\mathbf q\to0$ reproduces Eqs.~(408)--(410).
(iii)~$c\to\infty$ reproduces Eqs.~(446), (452) and (454), and with a tensor $\hat\epsilon^\infty$ the
standard non-analytic term.
(iv)~The transverse roots obey Huang's $c^2q^2=\omega^2\epsilon(\omega)$ and the Lyddane--Sachs--Teller
relation $\epsilon(0)/\epsilon_\infty=\omega_{\rm L}^2/\omega_\sigma^2$.
(v)~The Verdi--Giustino vertex [Giustino's Eq.~(180); SI with dimensionless $Z^*$, so
$Z^*_{\rm here}=|e|Z^*_{\rm Giustino}$ and $e^2/4\pi\varepsilon_0\to e^2$] equals the Fr\"ohlich coupling,
$\propto[1/\epsilon_\infty-1/\epsilon(0)]$, if and only if $\omega_{\rm LO}^2-\omega_{\rm TO}^2=\omega_p^2/
\epsilon_\infty$, whether $\epsilon(0)$ is taken from the Lyddane--Sachs--Teller relation or from
$\epsilon(0)=\epsilon_\infty+\omega_p^2/\omega_{\rm TO}^2$. This is the screened, unretarded L-mode shift of
(410c), which holds at every $\mathbf q$, i.e.\ for (410e) and (410f) alike. The unscreened shift of
Eq.~(410) satisfies it only for $\epsilon_\infty=1$, i.e.\ within its own electron-free model. The check
therefore tests the screening of the shift, not the order of limits.
Each of these checks, and the GaAs numbers above, was also verified numerically.

%======================================================================================
\section{The missing factor $e$ in the electrostatic Fr\"ohlich vertices (Sec.~7.2)}\label{sec:efactor}

\subsection{Location}
Sec.~7.2, pp.~93--99: the continuity equation on p.~94, Eqs.~(422)--(428), the displays on pp.~95, 96
and 98 and the sentence at the top of p.~99; also the 2D bubble on p.~101, and, for the same convention,
Eqs.~(234), (235), (237), (312) (pp.~47--48, 68) and one sentence in Sec.~6.

\subsection{Problem}
The connection between charge and current vertices, Eq.~(428), equates the charge vertex
$\rho_{n\mathbf k,m\mathbf k+\mathbf q}$ with $-\sum_\mu\hbar q_\mu j^\mu_{n\mathbf k,m\mathbf k+\mathbf q}/
(E_{n\mathbf k}-E_{m\mathbf k+\mathbf q})$. The former is a matrix element of the electron \emph{number}
density $\rho_{\rm e}=\Psi^+\Psi$ [Eqs.~(106), (239), (420), (426)]: $\rho_{\rm e}$ has the dimension of an
inverse volume, and the matrix element $\rho_{n\mathbf k,m\mathbf k+\mathbf q}$ is dimensionless. The latter
contains the charge $e$ through the current vertex (384). The factor is lost in the continuity equation
on p.~94, $\dot\rho_{\rm e}=-\nabla\mathbf j_{\rm e}$, which sets the rate of a number density equal to the
divergence of a charge current. For free electrons, $\phi_{\mathbf k}=e^{i\mathbf k\cdot\mathbf r}/\sqrt\Omega$
gives $\rho_{\mathbf k,\mathbf k+\mathbf q}=1$, $\mathbf j_{\mathbf k,\mathbf k+\mathbf q}=\frac{e\hbar}{2m_{\rm e}}
(2\mathbf k+\mathbf q)$ and $E_{\mathbf k}-E_{\mathbf k+\mathbf q}=-\frac{\hbar^2}{2m_{\rm e}}(2\mathbf k\cdot\mathbf q
+q^2)$, so $-\hbar\mathbf q\cdot\mathbf j/(E_{\mathbf k}-E_{\mathbf k+\mathbf q})=e$, not $1$. Both electrostatic
vertices obtained from (428), the bare one on p.~95 and the screened (boxed) one on p.~98, therefore
lack one factor $e$. As printed they have the dimension of $Z^*/(\Omega_{\rm uc}q)$ times a length, i.e.\
of an electrostatic potential, not of an energy. This is the $e^2$ ``missing'' in the margin note on
p.~96: one $e$ is inside $Z^*$, the other is the $e$ of (428). The propagator $D$ is not involved. The
same slip affects the transcription of Verdi and Giustino's vertex on p.~96 (with the manuscript's $Z^*$
it contains $e^3$) and the 2D electronic bubble on p.~101 (it needs $e^2$). The electrodynamic vertices
(414)--(417), (429)--(430), (447)--(448) and Eqs.~(431)--(446) are correct in this respect: there the
charge sits in $j^\mu$, $\Pi^{\rm e,0}_{jj}$ carries $e^2$, and so does $\Pi^{\rm e,0}_{\rho\rho}$ through
(350).

\subsection{Choice of fix}
Keep $\rho_{\rm e}=\Psi^+\Psi$ and $\rho_{n\mathbf k,m\mathbf k+\mathbf q}(\mathbf G)=\langle\phi_{n\mathbf k}|
e^{-i(\mathbf q+\mathbf G)\mathbf r}|\phi_{m\mathbf k+\mathbf q}\rangle$ as number densities, as they are defined
in (106), (239), (420), (426) and as the diamagnetic terms require: the $\mathbf A^2$ terms
$\frac{e^2}{2m_{\rm e}c^2}\int\rho_{\rm e}\mathbf A^2$ of (237), (270), (273), (312) and
$\Pi^{\rm e,dia}=-\frac{e^2}{mc}\rho^{\rm e}\delta(\mathbf r-\mathbf r')\delta_{\alpha\beta}$ of (279), (321).
Write $e\rho_{\rm e}$ wherever a charge density is meant. This leaves every definition, and every result
written with $j^\mu$, $j_{\alpha\sigma}$ or $W^{\rm e}$, unchanged. The alternative,
$\rho_{\rm e}\equiv e\Psi^+\Psi$ (the ``charge density operator'' of the text after (234)), would change
(239), (420), the $\mathbf A^2$ terms of (237), (270), (273), (312) and the diamagnetic terms of (279),
(321) ($e^2\to e$). Since (420) is substituted explicitly, it would also need the factor $e$ on the
left-hand sides of (423)--(425), of the two displays between (424) and (425) and of the display after
(427); only the continuity equation on p.~94 and (422) would stay as printed. And (428) would still need
the $e$ unless the vertex (106)/(426) were also redefined as $e\langle\phi|e^{-i(\mathbf q+\mathbf G)\mathbf r}|
\phi\rangle$, which would in turn change (105)--(107), (130)--(131) and the display on p.~96. Within
Sec.~7.2 the alternative therefore edits every equation edited here except the continuity equation and
(422); it adds eight edits elsewhere and saves only the Coulomb terms and the Sec.~6 sentence of
item~4 below.

\subsection{What changes, what does not}
\emph{Changed:} the display after (420) on p.~93 [(p93c$'$), a missing $\hbar$], the continuity equation on
p.~94 [(p94$'$)], Eq.~(422) [(422$'$)], the left-hand sides of (423)--(425), of the two displays between
(424) and (425) and of the display after (427), Eq.~(428) [(428$'$)], the vertices on pp.~95 and 98
[(p95$'$), (p98$'$)], the display and the margin note on p.~96 [(p96$'$) and new (428a)], one sentence on
p.~99 [new (448a)], the bubble on p.~101 [(p101$'$)], and, for the same convention, (234), (235), (237),
(312) [(234$'$), (235$'$), (237$'$), (312$'$)] and one sentence in Sec.~6. The displays on pp.~95--98 and
(237$'$) also carry the basis phase of Sec.~\ref{sec:phase}; their final forms are in
Sec.~\ref{sec:interplay}. \emph{Unchanged by the present correction} (several of them change by the basis
phase, Sec.~\ref{sec:phase}): (105)--(131) [Eq.~(131) already carries one explicit $e$], (239)--(240),
(321), (383)--(421), (426)--(427), (429)--(454), and (455)--(472) apart from the bubble on p.~101.

\subsection{Corrected equations and replacement prose}
\paragraph{1. Continuity and Eqs.~(422)--(428) (pp.~93--95).}
On p.~93 the Heisenberg evolution after (420) reads
\begin{equation*}
c_{n\mathbf k}(t)=e^{iH^{\rm e}_0t/\hbar}\,c_{n\mathbf k}\,e^{-iH^{\rm e}_0t/\hbar}.
\tag{p93c$'$}
\end{equation*}
On p.~94 replace ``So the equation for charge density is'' by ``So the equation for the electron (number)
density is''. Replace ``If we now use the continuity equation \ldots\ then we have (422)'' by: ``Since
$\rho_{\rm e}$ (420) is a number density and $\mathbf j_{\rm e}$ (419) a charge current density, the
continuity equation reads
\begin{equation*}
e\,\dot\rho_{\rm e}(\mathbf r,t)=-\nabla\cdot\mathbf j_{\rm e}(\mathbf r,t)
\tag{p94$'$}
\end{equation*}
and we have
\begin{equation*}
e\,[\rho_{\rm e}(\mathbf r,t),H^{\rm e}_0]_-=-i\hbar\,\nabla\cdot\mathbf j_{\rm e}(\mathbf r,t).\text{''}
\tag{422$'$}
\end{equation*}
Multiply the left-hand sides of (423), (424), (425), of the two unnumbered displays between (424) and
(425), and of the display after (427) by $e$. Eq.~(428) becomes
\begin{equation*}
e\,\rho_{n\mathbf k,m\mathbf k+\mathbf q}=-\sum_\mu\hbar q_\mu\,
\frac{j^\mu_{n\mathbf k,m\mathbf k+\mathbf q}}{E_{n\mathbf k}-E_{m\mathbf k+\mathbf q}} .
\tag{428$'$}
\end{equation*}
In the same passage: before (418), ``retarded limit $c\to\infty$'' should read ``non-retarded
(instantaneous) limit $c\to\infty$''; and before (417), ``reduce to standard bare Fr\"ohlich vertex
Eq.~(131)'' should read ``reduces, up to the overall sign $e/|e|$ [see (428a)] and the relabelling
$n\leftrightarrow m$, $\mathbf k\to\mathbf k+\mathbf q$, to the unscreened ($\hat\epsilon^\infty=1$),
$\mathbf G=0$ term of the polar vertex (131$'$), written in Gaussian units''. [Eq.~(131$'$) uses $e>0$,
with the electron charge written $-e$ in (123); from Sec.~4 on $e$ is the signed electron charge.]

\paragraph{2. The vertices on pp.~95, 96, 98 and the sentence on p.~99.}
With (428$'$) inserted in (418$'$) and (448$'$), both electrostatic vertices gain the factor $e$. Since
the basis phase of Sec.~\ref{sec:phase} changes the same displays, the final forms (p95$'$), (p96$'$),
(p98$'$), the relations (428a), (448a) to Verdi and Giustino's vertex, and the replacement prose for
pp.~95--99 are given once, in Sec.~\ref{sec:int-vertices}. The blue margin note on p.~96 can be deleted.

\paragraph{3. 2D electronic bubble (p.~101).}
The bubble that enters (435) and (461)--(470) is the charge--charge correlator [cf.\ (350)]:
\begin{equation*}
\Pi^{\rm e,0}_{\rho\rho}(\mathbf Q,\omega)=\frac{2e^2}{S}\sum_{\mathbf K\in{\rm SBZ}}\sum_{nm}
\frac{f_{n\mathbf K}-f_{m\mathbf K+\mathbf Q}}{\hbar\omega+E_{n\mathbf K}-E_{m\mathbf K+\mathbf Q}+i\eta}\,
\big|\rho_{n\mathbf K,m\mathbf K+\mathbf Q}\big|^2 .
\tag{p101$'$}
\end{equation*}
With it, (470) is dimensionless (with $v_Q=2\pi/Q$) and, for a 2D electron gas, gives Stern's
$\epsilon^{\rm e}_{\rm 2D}(Q,0)=1+2m^*e^2/(\hbar^2Q)$, valid for $Q\le2k_{\rm F}$ at $T\to0$ and without a
background dielectric constant. Eqs.~(461)--(462) and (469)--(472) are unchanged.

\paragraph{4. The same convention in the Hamiltonian (pp.~47--48, 68).}
Replace (234) and the sentence after it by
\begin{equation*}
\nabla\cdot\mathbf E=4\pi\big(e\rho_{\rm e}+\rho_{\rm ion}\big),
\tag{234$'$}
\end{equation*}
``where $\rho_{\rm e}=\Psi^+\Psi$ is the electron density operator (239), so that $e\rho_{\rm e}$ is the
electronic charge density operator'' [in (179)--(180) $\rho_{\rm e}$ denotes this charge density], and
(235) by
\begin{equation*}
H^0_{\rm Lph}=\frac12\int d\mathbf r\int d\mathbf r'\,
\frac{[e\rho_{\rm e}(\mathbf r)+\rho_{\rm ion}(\mathbf r)][e\rho_{\rm e}(\mathbf r')+\rho_{\rm ion}(\mathbf r')]}
{|\mathbf r-\mathbf r'|} .
\tag{235$'$}
\end{equation*}
The last term of (237) changes in the same way; (237) also changes by the basis phase, and its final
form (237$'$) is in Sec.~\ref{sec:int-hamiltonian}. Eq.~(312) becomes
\begin{equation*}
\begin{split}
V'={}&-\frac1c\int d\mathbf r\,\mathbf j^{\rm e}(\mathbf r)\mathbf A(\mathbf r)
+\frac{e^2}{2m_{\rm e}c^2}\int d\mathbf r\,\rho^{\rm e}(\mathbf r)\mathbf A^2(\mathbf r)\\
&-\frac1c\int d\mathbf r\,\mathbf j^{\rm ion}(\mathbf r)\mathbf A(\mathbf r)
+\frac{e^2}{2m_0c^2}\sum_{\alpha\beta}\int d\mathbf r\,A_\alpha(\mathbf r)\rho^{\rm ion}_{\alpha\beta}(\mathbf r)
A_\beta(\mathbf r)\\
&+\frac12\int d\mathbf r\int d\mathbf r'\,
\frac{[e\rho^{\rm e}(\mathbf r)+\rho^{\rm ion}(\mathbf r)][e\rho^{\rm e}(\mathbf r')+\rho^{\rm ion}(\mathbf r')]}
{|\mathbf r-\mathbf r'|};
\end{split}
\tag{312$'$}
\end{equation*}
its $\mathbf A^2$ terms stay as printed, and $\rho^{\rm ion}$, $\mathbf j^{\rm ion}$ and $\rho^{\rm ion}_{\alpha\beta}$ are
built from (314$'$), (272$'$) and (315$'$). In Sec.~6 add after (315): ``The electronic charge vertex is
therefore $e\rho^{\rm e}$, and $\Pi^{\rm e}_{\rho\rho}$, $\Pi^{\rm e}_{\rho j}$, $\Pi^{\rm e}_{j\rho}$ in
(320)--(350) are correlators of $e\rho^{\rm e}$, as the continuity equation (335) requires.'' These edits
change no displayed result.

\subsection{Justification}
Equation (428$'$) follows from the manuscript's own derivation once charge conservation is written for a
number density, (p94$'$). With free electrons it is checked in one line (see above), and it holds for any
local $H^{\rm e}_0$. Independently, the Coulomb term of (235$'$), $e\int d\mathbf r\,\rho_{\rm e}(\mathbf r)
\phi_{\rm ion}(\mathbf r)$, with $\phi_{\rm ion}$ the potential of $\rho_{\rm ion}=-\nabla\cdot\mathbf P$ and
$\mathbf P$ from (241$'$), gives as coefficient of $c^+_{m\mathbf k+\mathbf q}c_{n\mathbf k}
(a_{\mathbf q\sigma}+a^+_{-\mathbf q\sigma})$
\begin{equation*}
-\frac{4\pi ie}{\Omega_{\rm uc}}\sum_{\rm a}\sum_{\mu\nu}\sqrt{\frac{\hbar}{2Nm_{\rm a}\omega_\sigma(\mathbf q)}}\,
\frac{q_\mu Z^*_{{\rm a}\mu\nu}e_{{\rm a}\nu\sigma}(\mathbf q)}{q^2}\,
\langle\phi_{m\mathbf k+\mathbf q}|e^{i\mathbf q(\mathbf r-\mathbf r^0_{\rm a})}|\phi_{n\mathbf k}\rangle ,
\end{equation*}
i.e.\ Eq.~(131$'$) at $\hat\epsilon^\infty=1$, $\mathbf G=0$ in Gaussian units (with $n\leftrightarrow m$,
$\mathbf k\to\mathbf k+\mathbf q$). For electrons this is $g^0_{m n\sigma}$ of (p96$'$), and for either sign of
$e$ it is $-$(p95$'$). Checks: (i)~dimensions: $[eZ^*]$ is energy$\times$length, so (p95$'$), (p98$'$) are
energies, while the printed forms are energies divided by a charge; $v_Q\Pi^{\rm e,0}_{\rho\rho}$ of
(p101$'$) is dimensionless. (ii)~The longitudinal projection of (375),
$\hat{\mathbf q}\cdot D^0\cdot\hat{\mathbf q}=-4\pi c/\omega^2$, equals that of (374) $=$ (442) for every $c$,
and its factor $c$ cancels the $c^{-1}$ in the prefactor of (417). The longitudinal part of (417)
therefore does not depend on $c$, and (p95$'$) is that longitudinal part for every $c$. Likewise (p98$'$)
is the longitudinal part of (430) for an isotropic $\hat\epsilon^{\rm e}$
($\hat{\mathbf q}\cdot D^{\rm e}\cdot\hat{\mathbf q}=-4\pi c/(\omega^2\epsilon^{\rm e}_{\rm L})$). The transverse
parts vanish as $(\omega/cq)^2$. (iii)~With (p95$'$) and (p98$'$), $\sum g^{\rm dyn}\chi^0g^{\rm dyn,0}$
reproduces (413), and $\Omega^0_{\sigma\sigma}+\Delta\Omega^{\rm e}_{\sigma\sigma}$ reproduces (452), i.e.\ the
screened Gonze--Lee term; with the printed vertices both fail by $e^2$. (iv)~(428a), (448a) agree with
Giustino's Eq.~(180) (SI, dimensionless $Z^*$; $\mathbf G=0$, phase $e^{-i\mathbf q\bm{\tau}_\kappa}$
included); the numbers are in Sec.~\ref{sec:int-signs}. All of this was verified numerically, including
(428$'$) to $10^{-8}$ or better for Bloch states of a non-centrosymmetric local potential, the
$c$-independence of the longitudinal part in (ii) to $10^{-14}$, and Stern's result for (p101$'$).

%======================================================================================
\section{Interplay of the corrections}\label{sec:interplay}

Where two corrections touch the same equation, this section gives the single final form. The factor $e$
(Sec.~\ref{sec:efactor}) multiplies the prefactor; the basis phase (Sec.~\ref{sec:phase}) multiplies the
atom summand, with $e^{-i\mathbf q\mathbf r^0_{\rm a}}$ next to $e_{{\rm a}\nu\sigma}(\mathbf q)$ and
$e^{+i\mathbf q\mathbf r^0_{\rm a}}$ next to $e^*_{{\rm a}\nu\sigma}(\mathbf q)$. The two operations commute.

\subsection{The vertices on pp.~95--99 (factor $e$ and basis phase)}\label{sec:int-vertices}
\paragraph{p.~95.} Replace ``After using (428) in (418) we get'' and the display by ``After using (428$'$)
in (418$'$) we get''
\begin{equation*}
\lim_{c\to\infty}g^{\rm dyn,0}_{nm\sigma}(\mathbf k,\mathbf q,\omega)=
\frac{4\pi i\,e}{\Omega_{\rm uc}}\sum_{\rm a}\sum_{\mu\nu=xyz}
\sqrt{\frac{\hbar}{2Nm_{\rm a}\omega_\sigma(\mathbf q)}}\,
\frac{q_\mu Z^*_{{\rm a}\mu\nu}e_{{\rm a}\nu\sigma}(\mathbf q)\,e^{-i\mathbf q\mathbf r^0_{\rm a}}}{q^2}\,
\big[\rho_{n\mathbf k,m\mathbf k+\mathbf q}\big]^* .
\tag{p95$'$}
\end{equation*}

\paragraph{p.~96.} Replace the paragraph after (p95$'$), its display and the blue margin note by:
\begin{quote}
Let us write the long-range vertex of \textit{C.~Verdi and F.~Giustino, Phys.\ Rev.\ Lett.}\ \textbf{115},
176401 (2015), Eq.~4, with the basis phase restored as in Eq.~(180) of \textit{F.~Giustino, Rev.\ Mod.\
Phys.}\ \textbf{89}, 015003 (2017), for the moment without electronic polarizability
($\hat\epsilon^\infty=1$) and crystal local fields ($\mathbf G=0$), and in Gaussian units
($e^2/4\pi\epsilon_0\equiv e^2$). Their Born charges are dimensionless, ours carry the charge,
$Z^*_{{\rm a}\mu\nu}=|e|\,Z^{*\rm VG}_{{\rm a}\mu\nu}$, and their atomic positions $\bm{\tau}_\kappa$ are the
positions $\mathbf r^0_{\rm a}$ of our dipoles (314$'$). Their wave functions are normalized in the
Born--von K\'arm\'an supercell of $N_p=N$ cells, as ours are in $\Omega=N\Omega_{\rm uc}$, and their
$N_p^{-1/2}$ under the square root corresponds, as in (108), to $H_{\rm ep}=\sum g\,c^+c\,(a+a^+)$ without a
separate $N^{-1/2}$. In our notation their vertex reads
\begin{equation*}
\begin{split}
g^0_{mn\sigma}(\mathbf k,\mathbf q)={}&\frac{4\pi i\,|e|}{\Omega_{\rm uc}}\sum_{\rm a}\sum_{\mu\nu}
\sqrt{\frac{\hbar}{2Nm_{\rm a}\omega_\sigma(\mathbf q)}}\,
\frac{q_\mu Z^*_{{\rm a}\mu\nu}e_{{\rm a}\nu\sigma}(\mathbf q)}{q^2}\\
&\times\langle\phi_{m\mathbf k+\mathbf q}|e^{i\mathbf q(\mathbf r-\mathbf r^0_{\rm a})}|\phi_{n\mathbf k}\rangle .
\end{split}
\tag{p96$'$}
\end{equation*}
Since $\langle\phi_{m\mathbf k+\mathbf q}|e^{i\mathbf q(\mathbf r-\mathbf r^0_{\rm a})}|\phi_{n\mathbf k}\rangle
=e^{-i\mathbf q\mathbf r^0_{\rm a}}\big[\rho_{n\mathbf k,m\mathbf k+\mathbf q}\big]^*$, Eq.~(p95$'$) is
\begin{equation*}
\lim_{c\to\infty}g^{\rm dyn,0}_{nm\sigma}(\mathbf k,\mathbf q,\omega)=\frac{e}{|e|}\,g^0_{mn\sigma}(\mathbf k,\mathbf q)
=-g^0_{mn\sigma}(\mathbf k,\mathbf q).
\tag{428a}
\end{equation*}
The two agree, basis phase included, up to the sign of the electron charge. This overall sign is a
convention of the diagrammatic vertices (415)--(416), namely of the factor $i$ and of the order in which
the two vertices are written. It cannot be absorbed in $e_{{\rm a}\nu\sigma}(\mathbf q)$, because (p95$'$)
and (p96$'$) are linear in the same polarization vector. Only the product $g^{\rm dyn}\chi^0g^{\rm dyn,0}$
enters (413), and it is fixed: it does not depend on this convention or on the phase of
$e_{{\rm a}\nu\sigma}(\mathbf q)$, and together with $\Omega^0_{\sigma\sigma}$ it reproduces (452). The sign
therefore cancels in $\Delta\Omega^{\rm e}$ and in $|g|^2$. Where (p95$'$) is combined with vertices defined
as Hamiltonian matrix elements, e.g.\ a short-range part $g^{\rm S}$ as in Verdi and Giustino's
$g=g^{\rm S}+g^{\rm L}$, the Hamiltonian sign $g^0_{mn\sigma}=-\lim_{c\to\infty}g^{\rm dyn,0}_{nm\sigma}$ must
be used. Of Verdi and Giustino's $e^2$, one $e$ is contained in our $Z^*$ and the other is the $e$ of the
charge--current relation (428$'$); nothing is missing in $D$. The order of the band indices ($g_{mn}$
there, $g_{nm}$ here) is only a label.
\end{quote}

\paragraph{p.~98.} Replace ``and using (428) we get'' and the boxed display by ``and using (428$'$) we
get'' and the boxed display
\begin{equation*}
\lim_{c\to\infty}g^{\rm dyn}_{nm\sigma}(\mathbf k,\mathbf q,\omega)=
\frac{4\pi i\,e}{\Omega_{\rm uc}}\sum_{\rm a}\sum_{\mu\nu=x,y,z}
\sqrt{\frac{\hbar}{2Nm_{\rm a}\omega_\sigma(\mathbf q)}}\,
\frac{q_\mu Z^*_{{\rm a}\mu\nu}e^*_{{\rm a}\nu\sigma}(\mathbf q)\,e^{i\mathbf q\mathbf r^0_{\rm a}}}
{\mathbf q\cdot\hat\epsilon^{\rm e}(\mathbf q,\omega)\cdot\mathbf q}\,
\rho_{n\mathbf k,m\mathbf k+\mathbf q} .
\tag{p98$'$}
\end{equation*}

\paragraph{p.~99.} Replace ``This is again almost identical to \ldots\ electron-phonon vertex.'' by:
\begin{quote}
Compared with (p96$'$) and (p95$'$),
\begin{equation*}
\begin{aligned}
\lim_{c\to\infty}g^{\rm dyn}_{nm\sigma}(\mathbf k,\mathbf q,\omega)
&=-\frac{e}{|e|}\,\frac{q^2}{\mathbf q\cdot\hat\epsilon^{\rm e}(\mathbf q,\omega)\cdot\mathbf q}\,
\big[g^0_{mn\sigma}(\mathbf k,\mathbf q)\big]^*\\
&=-\frac{q^2}{\mathbf q\cdot\hat\epsilon^{\rm e}(\mathbf q,\omega)\cdot\mathbf q}\,
\Big[\lim_{c\to\infty}g^{\rm dyn,0}_{nm\sigma}(\mathbf k,\mathbf q,\omega)\Big]^* .
\end{aligned}
\tag{448a}
\end{equation*}
The second form follows from (415) and (416) alone: both vertices carry the factor $+i$ and the
propagator (442) is real, so the screened vertex is \emph{minus} the complex conjugate of the bare one
times $q^2/(\mathbf q\cdot\hat\epsilon^{\rm e}\cdot\mathbf q)$, whatever the sign convention for $e$. For real
$\hat\epsilon^{\rm e}$ and electrons, (448a) is the complex conjugate of the Verdi--Giustino vertex
($\mathbf G=0$), with $\hat\epsilon^\infty$ replaced by the dynamical electronic dielectric tensor
$\hat\epsilon^{\rm e}(\mathbf q,\omega)$; the basis phase becomes $e^{+i\mathbf q\mathbf r^0_{\rm a}}$ with
$e^*_{{\rm a}\nu\sigma}(\mathbf q)$. For $\hbar\omega$ well below the gap and $\mathbf q\to0$,
$\hat\epsilon^{\rm e}(\mathbf q,\omega)$ tends to the clamped-ion $\hat\epsilon^\infty$ (here the head of the
RPA dielectric matrix, without local fields and exchange--correlation), and the two vertices have the
same magnitude; for GaAs, $\sqrt N|g|\approx304$~meV at $q=0.1$~\AA$^{-1}$ from both (LO mode, the same
frequency $\omega_{\rm LO}$ under both square roots).
\end{quote}

\subsection{Why the signs are as stated}\label{sec:int-signs}
This subsection justifies (428a) and (448a); it is not meant for pasting. Insert $\lim_{c\to\infty}D^0$
(442) and $\lim_{c\to\infty}D^{\rm e}$ (446) into (415) and (416):
\begin{equation*}
\begin{aligned}
\lim_{c\to\infty}g^{\rm dyn,0}&=-\frac{4\pi ic}{\omega^2q^2}\,
(\mathbf q\cdot\mathbf J^{{\rm e}*})(\mathbf q\cdot\mathbf J_\sigma),\\
\lim_{c\to\infty}g^{\rm dyn}&=-\frac{4\pi ic}{\omega^2\,\mathbf q\cdot\hat\epsilon^{\rm e}\cdot\mathbf q}\,
(\mathbf q\cdot\mathbf J^*_\sigma)(\mathbf q\cdot\mathbf J^{\rm e}),
\end{aligned}
\end{equation*}
with $\mathbf J^{\rm e}$ the electronic vertex (414) and $\mathbf J_\sigma$ the ionic vertex (392), (385$'$).
Since $(-i)^*=+i$, $\lim g^{\rm dyn}=-[\lim g^{\rm dyn,0}]^*\,q^2/(\mathbf q\cdot\hat\epsilon^{\rm e}\cdot\mathbf q)$
for any complex $\hat\epsilon^{\rm e}$. The minus sign comes from the factor $+i$ that both (415) and (416)
carry; it does not involve $e$, and it holds with or without the basis phase, because the phase of the
ionic vertex is conjugated together with $e_{{\rm a}\nu\sigma}(\mathbf q)$. The overall sign $e/|e|$ of
(428a) is a different matter: it compares the diagrammatic vertex with a Hamiltonian matrix element. The
Coulomb route of Sec.~\ref{sec:efactor} (Justification) shows that the coefficient of
$c^+_{m\mathbf k+\mathbf q}c_{n\mathbf k}(a_{\mathbf q\sigma}+a^+_{-\mathbf q\sigma})$ in the Hamiltonian is
$-$(p95$'$) for either sign of $e$, and equals (p96$'$) for electrons. This is the convention of the
factorization (411)--(416). It cancels in (413), and it cannot be moved into the polarization vector.

With the dipoles at $\mathbf R^0_i$, as printed, the comparison with Verdi and Giustino would have to be
made at $\bm{\tau}_\kappa=0$. With the basis phase, $\bm{\tau}_\kappa=\mathbf r^0_{\rm a}$ in (p96$'$), in the
statement after (448a), and in check (iv) of Sec.~\ref{sec:efactor}. Both sides of (428a) and (448a) carry
the same phase, so the relations hold exactly as they would without it. The GaAs value 304.05~meV (LO mode, $\rho=1$,
$\mathbf q\to0$ eigenvectors; $Z^*=2.07\,|e|$, $\Omega_{\rm uc}=45.2$~\AA$^3$, reduced mass $36.1$~amu,
$\epsilon_\infty=10.9$, $\hbar\omega_{\rm LO}=36.2$~meV) is common to (p98$'$) with
$\hat\epsilon^{\rm e}\to\epsilon_\infty$ and to Giustino's Eq.~(180) evaluated in SI units. With finite-$q$
eigenvectors and the basis phase (force-constant model of Sec.~\ref{sec:phase}) the common value at
$q=0.1$~\AA$^{-1}$ is 303.0~meV along [111] and unchanged along [100]. Fr\"ohlich's vertex with
$1/\epsilon_\infty-1/\epsilon(0)$ from Lyddane--Sachs--Teller ($\hbar\omega_{\rm TO}=33.3$~meV) gives
319.8~meV; the ratio is $Z^*/Z^*_{\rm LST}=2.07/2.18$, a property of the input data. Numerically
verified: (428a), (448a) in both forms and their independence of the phase of $e_{{\rm a}\nu\sigma}$ for a
random three-atom cell with a complex symmetric $\hat\epsilon^{\rm e}$; the longitudinal parts of (415) with
(375) and of (416) with (410a) equal (p95$'$) and (p98$'$) at every $c$, with the full vertices converging
as $(\omega/cq)^2$; the cell-choice invariance (131a) of (p95$'$), (p96$'$), (p98$'$).

\subsection{E.410: (410c)--(410f) with the vertex (385$'$)}\label{sec:int-e410}
In Sec.~\ref{sec:e410}, (410b)--(410h) are written with $j_{\alpha\sigma}(\mathbf q)$, and the text states
that this is the phase-corrected vertex (385$'$). Written out, the two strengths in (410b), (410c) and
(410f) are
\begin{equation*}
\begin{aligned}
|\hat{\mathbf q}\cdot\mathbf j_\sigma|^2&=\Big|\sum_{{\rm a}\alpha\beta}\hat q_\alpha
\frac{Z^*_{{\rm a}\alpha\beta}}{\sqrt{m_{\rm a}}}e_{{\rm a}\beta\sigma}(\mathbf q)\,e^{-i\mathbf q\mathbf r^0_{\rm a}}\Big|^2,\\
|(1-\hat{\mathbf q}\hat{\mathbf q})\cdot\mathbf j_\sigma|^2&=\sum_\alpha\Big|\sum_{{\rm a}\beta}
\frac{Z^*_{{\rm a}\alpha\beta}}{\sqrt{m_{\rm a}}}e_{{\rm a}\beta\sigma}(\mathbf q)\,e^{-i\mathbf q\mathbf r^0_{\rm a}}\Big|^2
-|\hat{\mathbf q}\cdot\mathbf j_\sigma|^2,
\end{aligned}
\end{equation*}
and (410d) contains $j^*_{\alpha\sigma}(\mathbf q)$ with $e^{+i\mathbf q\mathbf r^0_{\rm a}}$. The final forms of
(410c)--(410f) are therefore those of Sec.~\ref{sec:e410}; no further display needs to be pasted. The
consequences are:
\begin{itemize}
\item At $\mathbf q\to0$ the phase is 1: (410$'$), (410e) and the GaAs values 55.6, 35.9 and 33.3~meV are
unchanged. So is the polariton region, where $q_c|\mathbf r^0_{\rm As}-\mathbf r^0_{\rm Ga}|\approx10^{-4}$.
\item At finite $\mathbf q$ both strengths are independent of the assignment of atoms to cells, and a rigid
acoustic translation, $e_{{\rm a}\beta\sigma}\propto\sqrt{m_{\rm a}/M}\,u_\beta e^{i\mathbf q\mathbf r^0_{\rm a}}$,
has $\mathbf j_\sigma\propto\sum_{\rm a}\hat Z^*_{\rm a}\mathbf u=0$ by the acoustic sum rule: (410c) gives the
acoustic branches no dipolar shift beyond the internal-strain term of (131b). With the phase-less vertex
it would: for LA[111] in GaAs the longitudinal term of (410c) alone gives $\sqrt{\Delta\omega^2}/q=2.5$ or
5.9~km/s depending on the cell, comparable to the sound velocity 5.4~km/s.
\item For $cq\gg\sqrt{\epsilon_\infty}\,\omega$, (410f) with (385$'$) is (454) with (385$'$), and (410d),
written in the atom basis, is the $\mathbf G=0$ term of (454a):
\begin{equation*}
\begin{split}
&\sum_{\sigma\sigma'}e_{{\rm a}\mu\sigma}(\mathbf q)\Big[\omega^2_\sigma\delta_{\sigma\sigma'}
+\frac{4\pi}{\Omega_{\rm uc}}\frac{(\hat{\mathbf q}\cdot\mathbf j_\sigma)^*(\hat{\mathbf q}\cdot\mathbf j_{\sigma'})}
{\hat{\mathbf q}\cdot\hat\epsilon^\infty\cdot\hat{\mathbf q}}\Big]e^*_{{\rm b}\nu\sigma'}(\mathbf q)\\
&\qquad=\tilde\Omega^2_{0,{\rm a}\mu,{\rm b}\nu}(\mathbf q)
+\frac{4\pi}{\Omega_{\rm uc}\sqrt{m_{\rm a}m_{\rm b}}}\,
\frac{(\mathbf q\hat Z^*_{\rm a})_\mu(\mathbf q\hat Z^*_{\rm b})_\nu}{\mathbf q\cdot\hat\epsilon^\infty\cdot\mathbf q}\,
e^{i\mathbf q(\mathbf r^0_{\rm a}-\mathbf r^0_{\rm b})}.
\end{split}
\end{equation*}
This was verified numerically for a random three-atom cell with an anisotropic $\hat\epsilon^\infty$; the
deviation decreases as $(\omega/cq)^2$.
\end{itemize}

\subsection{The interaction Hamiltonian (237)}\label{sec:int-hamiltonian}
Eq.~(237) is touched by both corrections: by the rule of Sec.~\ref{sec:phase} ($A_\alpha$ at
$\mathbf R^0_{ia}$) and by the factor-$e$ convention of Sec.~\ref{sec:efactor} ($e\rho_{\rm e}$ in the
Coulomb term). Its final form is
\begin{equation*}
\begin{split}
V'={}&-\frac1c\int d\mathbf r\,\mathbf j_{\rm e}(\mathbf r)\mathbf A(\mathbf r)
+\frac{e^2}{2m_{\rm e}c^2}\int d\mathbf r\,\rho_{\rm e}(\mathbf r)\mathbf A^2(\mathbf r)
-\frac1c\sum_{ia\alpha\mu}\frac{1}{m_{\rm a}}Z^*_{{\rm a}\alpha\mu}\pi_{ia\mu}A_\alpha(\mathbf R^0_{ia})\\
&+\frac{1}{2c^2}\sum_{ia\mu\alpha\beta}\frac{1}{m_{\rm a}}Z^*_{{\rm a}\alpha\mu}Z^*_{{\rm a}\beta\mu}
A_\alpha(\mathbf R^0_{ia})A_\beta(\mathbf R^0_{ia})\\
&+\frac12\int d\mathbf r\int d\mathbf r'\,
\frac{[e\rho_{\rm e}(\mathbf r)+\rho_{\rm ion}(\mathbf r)][e\rho_{\rm e}(\mathbf r')+\rho_{\rm ion}(\mathbf r')]}
{|\mathbf r-\mathbf r'|},
\end{split}
\tag{237$'$}
\end{equation*}
with $\rho_{\rm ion}=-\nabla\cdot\mathbf P$ and $\mathbf P$ from (241$'$). The $\mathbf A^2$ term keeps
$\rho_{\rm e}=\Psi^+\Psi$. (235$'$) and (312$'$) are touched by the factor-$e$ convention only; their
$\rho_{\rm ion}$ and $\rho^{\rm ion}$ are built from the phase-corrected $\mathbf P$ of (241$'$), (314$'$), so
no further change is needed. (186$'$) and (270) carry no Coulomb term and change by the basis phase only.

\subsection{Order of the new material on p.~99}\label{sec:int-p99}
Both the E.410 correction and the basis phase add text after Eq.~(454). Including the replacement at the
top of p.~99, p.~99 then reads, in this order:
\begin{enumerate}
\item the replacement of ``This is again almost identical \ldots'' with (448a)
(Sec.~\ref{sec:int-vertices});
\item the title of Sec.~7.3, ``Phonon's self-energy in electrostatic limit and the LO--TO splitting''
(Sec.~\ref{sec:e410}, item~5);
\item Eqs.~(449)--(454), unchanged in form, with the vertex (385$'$);
\item the sentence after (454) on Sohier \textit{et al.}\ (possibly reworded, Sec.~\ref{sec:minor}, p.~99);
\item the E.410 paragraph ``With $W^{\rm e}(\mathbf q,0)=4\pi/(\mathbf q\cdot\hat\epsilon^\infty\cdot\mathbf q)$,
Eq.~(454) is the standard non-analytic term \ldots'' (Sec.~\ref{sec:e410}, item~5);
\item the paragraph with (454a): ``Eq.~(454) keeps its form with (385$'$). For an insulator with
$W^{\rm e}_{\mathbf G\mathbf G'}=\ldots$, (399a) gives the matrix, $\mathbf G$-resolved form of (454)''
followed by (454a) and the three sentences after it (Sec.~\ref{sec:phase-678}).
\end{enumerate}
The E.410 paragraph comes first because it interprets (454) itself, the macroscopic $\mathbf G=0$ term,
and links it back to Sec.~7.1. The paragraph with (454a) generalizes (454) to $\mathbf G\neq0$; it ends
Sec.~7.3 with the remark on regularization that leads to the open point on point-dipole self-energies
(Sec.~\ref{sec:minor}, p.~48). Sec.~8 then starts on p.~100 as printed.

\subsection{Index of corrected and new equations, in manuscript order}\label{sec:int-index}
Primes: corrected (replace the printed equation or display). Letters: new. B, L, F mark the
corrections of Secs.~\ref{sec:phase} (basis phase), \ref{sec:e410} (E.410) and \ref{sec:efactor} (factor
$e$); B+F marks a merged form from Sec.~\ref{sec:interplay}.
\begin{itemize}
\item Sec.~3, pp.~24--25: (121$'$), (122$'$), (123$'$), (124$'$), (131$'$), (131a), (131b) [B].
\item Sec.~4, pp.~26--49: (134$'$), (143$'$), (167$'$), (183$'$), (186$'$) [B]; (234$'$), (235$'$) [F];
(237$'$) [B+F]; (241$'$, 314$'$) [B].
\item Sec.~5, pp.~56--67: (271$'$, 315$'$), (272$'$), (274$'$), (290$'$), (293$'$), (295$'$), (296$'$),
(297$'$), (299$'$), (300$'$), (302$'$), (304$'$), (305$'$), (307$'$), (308$'$), (309$'$) [B].
\item Sec.~6, pp.~68--86: (312$'$) [F]; (352$'$), (354a), (385$'$) [B].
\item Sec.~7, pp.~89--99: (399a) [B]; (410$'$), (410a)--(410h) [L, with (385$'$)]; (p93a$'$), (p93b$'$),
(418$'$) [B]; (p93c$'$), (p94$'$), (422$'$), (428$'$) [F]; (p95$'$), (p96$'$), (428a) [B+F]; (430$'$),
(447$'$), (448$'$) [B]; (p98$'$), (448a) [B+F]; (454a) [B].
\item Sec.~8, p.~101: (p101$'$) [F].
\end{itemize}
Corrected by a stated rule and not displayed: (141), (142), (270), (286), (289), (291), (303), (306),
(310), (311), (354) and the unnumbered displays on pp.~38, 39, 61, 65, 76--77 [B]; the left-hand sides of
(423)--(425), of the two displays between (424) and (425) and of the display after (427) [F]. Symbol
slips fixed in place: (406), (121) (label, included in (121$'$)).

%======================================================================================
\section{Minor errata}\label{sec:minor}
The following were collected while the primer was checked against the PDF. Items already covered above
are not repeated: the label of (121), the citation after (131), the Heisenberg evolution on p.~93,
``retarded limit'' before (418), the comparison with (131) before (417), the photon/phonon slips before
(406)--(407), and $\rho_{\rm e}$ in (234)--(237), (312).
\begin{itemize}\setlength{\itemsep}{0pt}
\item p.~1 (before E.3), p.~18: $|\mathbf Q_{\rm a}|\ll|\mathbf r^0_{\rm a}|$ depends on the origin; compare with the interatomic distance, write $\mathbf Q_{ia}$.
\item p.~2, E.4: state the on-site block $k_{ia\mu,ia\nu}=-\sum_{jb\neq ia}k_{ia\mu,jb\nu}$ (acoustic sum rule).
\item p.~4, E.15 (E.35, pp.~8--9): use $e^{-i\mathbf k\bm{\delta}}$ (Giustino's Eq.~14); required by Sec.~\ref{sec:phase}.
\item p.~10, before E.51: ``(46)'' $\to$ ``(50)'' (the step uses $[a,a^+]=1$).
\item p.~14, before E.76: $\omega^2_\alpha(-\mathbf k)=\omega^2_\alpha(\mathbf k)$ follows from time reversal, not from isotropy.
\item p.~14, after E.78: exclude the three $\mathbf k=0$ translations ($3Ns-3$ oscillators); $1/\sqrt\omega$ diverges in E.72, E.83.
\item p.~20, after E.108: $g^0$ is the bare linear coupling, not ``the first non-adiabatic correction''.
\item p.~21: the bare phonon propagator $D^0$ should be $F^0$ ($D$ is the photon propagator from Sec.~4.5 on).
\item p.~21: DFPT force constants of polar insulators have an $R^{-3}$ dipole tail; call the one-cell $k^{\mathcal S}$ a model.
\item p.~22, E.110: define $Z^*$ at zero macroscopic field (transverse charge), as E.120 assumes.
\item p.~23: announce at the start of Sec.~4 the switch from SI to Gaussian units.
\item p.~24 (E.123), pp.~26--27: state at E.134 the change from $-e$ ($e>0$) to the signed charge $e$.
\item p.~26, E.134: $k$ must be the short-range $k^{\mathcal S}$ of p.~21; the scheme covers the polar channel only.
\item pp.~26, 56, 57, 59, 70, 83: steps deferred to ``Lecture~9''/``Chapter~9''; include them or cite a source.
\item p.~27: $\mathbf a\times(\mathbf b\times\mathbf c)=\mathbf b(\mathbf a\cdot\mathbf c)-(\mathbf a\cdot\mathbf b)\mathbf c$, with $\nabla_e$ acting on $\mathbf A$ only.
\item p.~30: $k^\mu_{\rm ba}\to k^\mu_{\rm bb}$; use the index order of E.144--E.145 throughout.
\item pp.~30--31: $k_{\rm ab}=-k_{\rm aa}$ etc.\ hold for the lattice-summed constants (sum rule), not for each pair.
\item p.~32, before E.150: $Z^*_{\rm a}=-Z^*_{\rm b}$ is exact for a diatomic crystal (charge neutrality).
\item p.~32, after E.151: the Lorentz field $\mathbf E+\frac{4\pi}{3}\mathbf P$ holds for cubic sites only.
\item pp.~33--34, E.158--E.160: the local-field model is illustrative; only $\tilde\omega^2_{\rm LO}-\tilde\omega^2_{\rm TO}=\omega_p^2/\epsilon_\infty$ is general.
\item p.~34, before E.162: the zeros of $\epsilon$ give the LO modes only; the TO modes come from $\epsilon\to\infty$.
\item p.~35 (E.166), p.~48 (E.237--E.238): kinetic and paramagnetic currents need two symbols.
\item pp.~41--42, E.196--E.201: ``orthogonal'', not unitary; left-handed triad; $\mathbf e_3(\mathbf q)=\pm\hat{\mathbf q}$.
\item pp.~43--44, 52--55: photon operators $b_{\mathbf q\sigma}$ throughout (not $a_{\mathbf q\alpha}$).
\item p.~45, below E.219: phonons $a_{\mathbf q\alpha}$, photons $b_{\mathbf q\sigma}$ (labels swapped in the text).
\item p.~48, E.235: normal-order the electron term, drop on-site dipole self-energies; (454a) needs this too.
\item p.~50: Matsubara denominator $i\hbar\omega_n+\xi_{\mathbf k}-\xi_{\mathbf k+\mathbf q}$.
\item p.~50: state the scope of ``set $n_{\mathbf q\sigma}=0$ in boson lines'' (fails inside electron self-energies).
\item p.~53, E.260; p.~54: one factor $2\omega/(\omega^2-\omega^2_{\mathbf q}+i\eta)$, not two; $D^0\to D^{0t}$ in the middle step.
\item p.~54 (E.263), pp.~80, 83: explain $D^{0l}=0$ (interaction picture) next to E.374 (instantaneous Coulomb).
\item p.~57, E.276: start the series at $n=0$; write $V^t_I$ throughout.
\item pp.~58--60, 69 (E.279, E.284, E.321): $m\to m_{\rm e}$; ground-state average of $\rho^{\rm e}$ in $\Pi^{\rm e,dia}$.
\item p.~60, E.285; p.~86: a sheet needs the $z$-resolved Dyson equation; $l$ is the superlattice period.
\item p.~60, E.286--E.287: $\alpha,\beta$ serve as Cartesian and branch indices; use $\sigma$ for branches.
\item p.~64, before E.300: ``instead of (354)'' $\to$ ``instead of (291)''.
\item p.~65, E.301: $\Pi^{\rm dia}=-\Pi^{\rm para}(0)$ follows from gauge invariance; ions carry the bound charge $-\nabla\cdot\mathbf P$.
\item pp.~70, 75, 77: the connection formulas follow from charge conservation (Ward identities); state when.
\item p.~74: $\Pi_{j\rho}\to\Pi_{\rho j}$; $\int d\mathbf r\,\rho\,\phi^{\rm ext}$ (stray $V$).
\item pp.~78--79, 83: label the scalar slot 0, not 1; $V\to v$.
\item p.~82, Fig.~10(c): $j^{\rm ie}\to j^{\rm e}$.
\item p.~85, E.380: say it is the $\mathbf G=\mathbf G'=0$ element (see the p.~85 text of Sec.~\ref{sec:phase}).
\item p.~87 (E.391), p.~100 (E.455): write $\Pi^{\rm ion,0}_{jj}$ for the bare correlator.
\item p.~91, Fig.~14: draw $D^{\rm e}$ in the blue of Fig.~13.
\item p.~92 (E.414): $\chi^0$ is not the Lindhard function; only $J\chi^0J^*$ is regular for degenerate states.
\item p.~97, E.440: $D^{lt},D^{tl}\sim c^{-1}$, not $c^0$.
\item p.~99, before E.454: static $W^{\rm e}$ only for gapped systems; keep E.453 dynamical in doped ones.
\item p.~99, after E.454; p.~102: the comparison with Sohier \textit{et al.}\ is not carried out; state agreement in form.
\item pp.~60, 100--102: define SBZ (surface Brillouin zone) and the thickness $l$; ``the the'' (p.~102).
\item p.~101, before E.463: E.463--E.472 hold for in-plane modes; out-of-plane $Z^*_{zz}$ gives a $Q$-independent shift.
\item Spelling: ``communication relations'' (p.~14) should be ``commutation relations''; other slips (e.g.\ ``equlibrium'', ``Moreower'', ``Faynman'', ``Poison's'') are harmless.
\end{itemize}

\end{document}
